\documentclass[acmsmall,review,screen,nonacm]{acmart}

% Preserve the repository's single-column ACM review layout.
% TODO: Replace author metadata and add venue metadata and CCS concepts.
\usepackage{amsmath,booktabs,tabularx,array}
\usepackage{listings}
\usepackage{placeins}
\usepackage{float}
\usepackage{microtype}
\usepackage{needspace}
\setcounter{secnumdepth}{3}
\setcounter{tocdepth}{3}
\lstdefinelanguage{TEPL}{
  morekeywords={dialect,op,attrs,axes,semantics,dtype,bind,yield,let,rule,where,derive,from,import,as,use,auto,tensor,index},
  sensitive=true,morecomment=[l]{//},morestring=[b]"
}
\lstset{basicstyle=\normalfont\ttfamily\small,keywordstyle=\bfseries,
  commentstyle=\normalfont\ttfamily\color{black!60},
  columns=fullflexible,keepspaces=true,showstringspaces=false,
  breaklines=true,breakatwhitespace=true,frame=single,
  framerule=0.3pt,rulecolor=\color{black!25},xleftmargin=4pt,xrightmargin=4pt,
  aboveskip=8pt,belowskip=8pt,captionpos=b}
\newcommand{\code}[1]{{\upshape\ttfamily\detokenize{#1}}}
\newcommand{\todo}[1]{\textcolor{black!75}{\textbf{[TODO:} #1\textbf{]}}}
\newcommand{\catseq}{\mathbin{+\!\!+}}
\newcommand{\gath}{\operatorname{gather}}
\newcommand{\ocomp}{\operatorname{ordered\_complement}}
\newcommand{\pos}{\operatorname{positions}}
\newcommand{\mpo}{\operatorname{matchPositions}}
\newcommand{\Axes}{\operatorname{Axes}}
\newcommand{\Keep}{\operatorname{KeepAxes}}
\newcommand{\Isel}{I_{\mathrm{sel}}}
\newcommand{\Irem}{I_{\mathrm{rem}}}
\newcommand{\Jsel}{J_{\mathrm{sel}}}
\newcommand{\Jrem}{J_{\mathrm{rem}}}
\newcolumntype{Y}{>{\raggedright\arraybackslash}X}
% Float barriers keep every figure in its specified subsection.
% Small artwork uses 50% of text width; large artwork uses 100%.
\floatstyle{plaintop}
\newfloat{paperalgorithm}{htbp}{loa}
\floatstyle{plain}
\newfloat{codeblock}{htbp}{loc}
\floatname{paperalgorithm}{Algorithm}
\newcommand{\paperfigure}[5]{%
  \begin{figure}[!htbp]\centering
  \captionsetup{width=#1\linewidth}
  \includegraphics[width=#1\linewidth]{#2}
  \caption{#3}\label{#4}\Description{#5}
  \end{figure}}


% Put 50%-width artwork beside prose occupying 46% of text width. Measuring both columns keeps
% each pair on one page and prevents text or display math from crossing artwork.
\newsavebox{\figuretextbody}
\newsavebox{\figuretextart}
\newlength{\figuretextheight}
\newenvironment{figuretext}[4]{%
  \FloatBarrier
  \def\figuretextfile{#1}\def\figuretextcaption{#2}%
  \def\figuretextlabel{#3}\def\figuretextdescription{#4}%
  \begin{lrbox}{\figuretextbody}%
  \begin{minipage}[t]{0.46\textwidth}\vspace{0pt}%
  \emergencystretch=1em\noindent
}{%
  \end{minipage}\end{lrbox}%
  \begin{lrbox}{\figuretextart}%
  \begin{minipage}[t]{0.5\textwidth}\vspace{0pt}\centering
    \captionsetup{type=figure,hypcap=false,justification=raggedright,singlelinecheck=false}%
    \includegraphics[width=\linewidth]{\figuretextfile}%
    \caption{\figuretextcaption}\label{\figuretextlabel}%
    \Description{\figuretextdescription}%
  \end{minipage}\end{lrbox}%
  \setlength{\figuretextheight}{\dimexpr\ht\figuretextbody+\dp\figuretextbody\relax}%
  \ifdim\dimexpr\ht\figuretextart+\dp\figuretextart\relax>\figuretextheight
    \setlength{\figuretextheight}{\dimexpr\ht\figuretextart+\dp\figuretextart\relax}%
  \fi
  \typeout{SIDEFIG \figuretextlabel: \the\figuretextheight}%
  \par\addvspace{\intextsep}%
  \Needspace{\dimexpr\figuretextheight+\baselineskip\relax}%
  \noindent\usebox{\figuretextbody}\hfill\usebox{\figuretextart}\par
  \addvspace{\intextsep}%
}

\title[TEPL]{TEPL: Synthesizing Attribute Programs from Tensor Rewrite Semantics}
\author{Author Name}
\affiliation{\institution{Institution}\city{City}\country{Country}}

\begin{document}
\raggedbottom
\begin{abstract}
Tensor rewrite rules are concise in algebraic form, but their implementations
require substantial attribute-transformation logic. Symbolic patterns express
computation structure; matching alone does not derive the permutations,
contraction dimensions, and other attributes needed by the replacement.
TEPL (Tensor Equality Pattern Language) lets users describe tensor rewrites,
explicit conditions, and attribute changes together. When replacement
attributes are omitted, its compiler infers them from operation semantics and
checks supported rewrites for legality and equivalent results. The central
insight is that coordinate correspondences used to establish semantic
equivalence also provide constructive witnesses for attribute synthesis.
TEPL combines arbitrary-length axis sequences, compositional ordered mappings,
lazy partition refinement, and algebraic normalization to construct reusable
functions of source metadata. The intended correctness contract requires
computable replacement attributes and preservation of legality and tensor
semantics under explicit preconditions and a selected arithmetic policy;
this draft states that contract without claiming an established soundness
result for the implementation. We describe the language, synthesis procedure,
and implementation, and plan an evaluation of authoring effort, remaining
attribute expressions and conditions, semantic-definition cost,
production-derived coverage, and synthesis performance. Measurements and
results remain to be collected.
\end{abstract}
\keywords{tensor rewrites, attribute synthesis, symbolic axis mapping, rewrite languages}
\maketitle

\FloatBarrier
\FloatBarrier
\section{Introduction}
\label{sec:introduction}
\FloatBarrier
\subsection{Motivation}
Algebraic tensor rewriting is used to optimize tensor computations by replacing
one computation graph with an equivalent graph. The equation describing a
rewrite can be short even when its implementation is complex. For example,
reassociating a product from $(XY)Z$ to $X(YZ)$ changes only the nesting of two
dots in an algebraic expression. For generalized tensor dots, however, an
implementation must also reconstruct which dimensions are batched or
contracted, preserve their pairings, and arrange the output dimensions in the
required order. XLA's \code{AssociativeReorderNestedDot} routine illustrates
these tasks: it constructs new dimension descriptors and, when needed, an
output permutation, alongside applicability and cost checks~\cite{xla-simplifier}.
This observation concerns the kinds of work required, rather than a measured
line-count comparison. \todo{Classify the pinned XLA routine's attribute logic
separately from legality, profitability, and graph construction before
reporting its size.}

Symbolic rewriting frameworks such as egg support concise operation
patterns~\cite{egg}. They can represent operation attributes, but their pattern
matching does not itself derive attribute transformations from tensor
semantics. For a tensor rewrite, the user must supply the additional
expressions or custom host-language logic; egg's \code{Applier} interface, for
example, permits application code to inspect analysis data and construct a
replacement~\cite{egg-applier}. Merely matching the operations and attributes
on the left-hand side (LHS) does not determine the attributes needed on the
right-hand side (RHS).

Manual attribute transformation is a major source of this complexity, although
legality checks, numerical policy, and profitability also contribute. Cascaded
operations change axis positions, operand pairings, and output order. Equal
dimension sizes do not establish that two dimensions denote the same logical
coordinate. Moreover, the transformation should work for tensors of arbitrary
rank; working out a few concrete matrix or fixed-rank examples does not
establish the general relationship.

Our research question is therefore: \emph{What if we could design a language
that makes tensor rewriting as easy and intuitive as writing its algebraic
form, while leaving complex attribute transformation to the compiler?}
TEPL pursues this goal through operation semantics, symbolic axis mapping,
and automatic RHS attribute inference under explicit conditions.
Figure~\ref{fig:why-tepl} compares the authoring choices for the same
reassociation. A bare algebraic pattern supplies the replacement graph but
leaves both replacement descriptors unspecified. Custom implementation code
supplies them through handwritten dimension logic. TEPL supports both
user-written descriptor expressions and compiler inference from reusable
operation semantics. The example is unbatched and presupposes arithmetic
that permits reassociation; it does not establish that every generalized-dot
reassociation is valid.

\paperfigure{1}{figures/Why TEPL_ From Algebraic Rules to Practical Tensor Rewrites.pdf}{TEPL combines concise tensor rewrite patterns with explicit or inferred replacement attributes.}{fig:why-tepl}{Four panels compare an algebraic pattern, custom implementation logic,
explicit TEPL attribute expressions, and inferred TEPL attributes for the
same nested-dot replacement graph.}

\FloatBarrier
\subsection{Central Thesis}
TEPL describes tensor rewrites, their conditions, and their attribute changes
together. Users can write attribute changes directly or ask the compiler to
infer them from operation semantics. The compiler checks supported semantic
rewrites for validity and equivalent results.

The language lets users express algebraic patterns, applicability conditions,
and attribute changes in one place. It supports user-defined dialects, each
with its own operations, attributes, and semantics. An operation is defined
once within a dialect and then imported and reused across rules. Tensor
operations, graph patterns, attribute bindings, metadata conditions, and
replacement-attribute expressions are represented directly in the language.

The central technical advance pursued here is inference of omitted RHS
attributes. Users may also request this inference explicitly with
\code{auto}. Symbolic axis mapping constructs and checks reusable attribute
programs for arbitrary-rank tensors within the supported semantics. The
compiler reasons from operation legality and the user's stated conditions;
when those assumptions do not establish correctness, the semantic rule is
rejected as unproved. Inference does not silently add conditions to make a
proposed rewrite succeed.

\FloatBarrier
\subsection{Proposed Contributions}
This paper develops three connected contributions:
\begin{enumerate}
\item \textbf{A language for tensor rewrites.} TEPL supports user-defined
dialects and operations and expresses rewrite patterns, applicability
conditions, and attribute changes together. Users can capture and reuse
descriptors or compute replacement attributes directly in the language.
\item \textbf{Automatic attribute inference.} The compiler synthesizes omitted
RHS attributes for supported rewrites as reusable functions of LHS metadata,
without fixing tensor rank. Its Axes IR builds on prior aggregated-axis
representations and adds explicit ordered interfaces, operand mappings,
and value semantics for this synthesis task.
\item \textbf{Symbolic axis mapping for composed operations.} The method tracks
logical-axis correspondence across composed operations, splits and composes
arbitrary-length groups while preserving coordinate order and pairings,
and reuses established producer--consumer mappings for algebraic matching
and attribute construction.
\end{enumerate}
These are the design and mechanism claims of this draft. The correctness
contract is stated as an intended guarantee, and the practical benefits are
questions for the planned evaluation, rather than established experimental
results.

\FloatBarrier
\FloatBarrier
\section{Background}
\label{sec:background}
TensorRight verifies tensor graph rewrites for tensors of arbitrary rank and
size~\cite{tensorright}. Its aggregated axes represent groups containing an
arbitrary number of named axes, allowing one rewrite specification to describe
tensors with different ranks. Grouping related axes makes it possible to reason
about a family of tensors without enumerating each dimension separately.

This representation provides the background for TEPL's symbolic axis groups.
TEPL retains the ability to describe arbitrary-length groups but represents
them as ordered sequences and tracks their positions in tensor interfaces.
The order matters for constructing concrete operation attributes: a pair of
contraction lists specifies not only which axes are reduced, but which axes
are paired and in what list order. The distinction between a logical axis,
its extent, and its position is central to the method developed in
Section~\ref{sec:method}. Axis grouping itself is prior work; the contribution
pursued here is its use with ordered, compositional correspondence to
construct reusable attribute programs.

\FloatBarrier
\FloatBarrier
\section{High-Level Language Design}
\label{sec:language}
\FloatBarrier
\subsection{Dialects, Operations, and Attributes}
\label{sec:dialects}
A dialect groups user-defined operations with their operand signatures,
attribute schemas, and semantics. Users define an operation once and import
it into multiple rules; different dialects can provide different operation
vocabularies. Attributes configure an operation independently of its tensor
operands. For example, a transpose carries a permutation list. The following
excerpt declares its signature and attribute in a dialect file,
\code{tensor.tepl}:
\begin{codeblock}[H]
\begin{lstlisting}[language=TEPL]
dialect TensorLang {
  op transpose(input: tensor) -> tensor {
    attrs {
      permutation: index[];
    }
  }
}
\end{lstlisting}
\end{codeblock}
This excerpt introduces the schema; it omits the semantic blocks. In a
semantic definition, \code{dtype} specifies element-type behavior,
\code{axes} declares coordinate bindings, legality, and output order, and
\code{semantics} gives the value at those coordinates. For transpose,
\code{dtype} preserves the element type. Transpose changes output coordinate
order while preserving the coordinate-indexed input value, so a value clause
\code{yield input} does not imply an unchanged tensor interface. We defer
the representation and composition of these blocks to Axes IR in
Section~\ref{sec:axes-ir}.

\FloatBarrier
\subsection{Imports and Rewrite Rules}
A rule imports a dialect and selects the operations it uses:
\begin{codeblock}[H]
\begin{lstlisting}[language=TEPL]
from "tensor.tepl" import TensorLang as t;
use t::{transpose};
\end{lstlisting}
\end{codeblock}
Rewrites are written as algebraic equations with an LHS pattern followed by an
RHS replacement. Nested operations denote producer--consumer edges, while
tensor captures identify matched values or subcomputations. Repeated captures
refer to the same matched value. This application graph is distinct from the
semantic graph used to reason about tensor elements and their coordinates.

Our running pattern is
\code{(transpose[@outer] (transpose[@inner] X))}.
Its intended replacement is a single transpose of \code{X}; the remaining
task is to construct that transpose's permutation. Tensor declarations may
also constrain shape and dtype, but equal dimension sizes alone do not
establish coordinate identity.

\FloatBarrier
\subsection{Attribute Binding and Reuse}
\label{sec:binding}
The annotation \code{op[@name]} binds an operation's attribute descriptor on
the LHS. If LHS attributes are omitted, the matcher captures each operation's
descriptor independently and anonymously. A name is needed when an expression
or an RHS operation refers to that descriptor. Captures are available in
\code{where} conditions, \code{derive} expressions, and RHS operations.
Field access such as \code{@outer.permutation} selects an individual
attribute where the expression context supports it.

\begin{figuretext}{figures/Attribute Binding and Derivation.pdf}{Captured operation descriptors can be read and transformed to configure replacement operations.}{fig:attribute-binding}{An input graph with transpose and dot is replaced by a dot. Colored descriptor
bindings connect the rule to a table with two captured rows and one derived row.}
Reusing \code{@name} on a compatible RHS operation reuses its descriptor;
a derived descriptor can configure an RHS operation in the same way.
Compatibility checks the attribute structure, rather than the mathematical
correctness of that choice. In the transpose-composition example, binding
\code{@inner} and \code{@outer} makes both permutations available, but copying
either one alone generally does not give the replacement permutation.

Figure~\ref{fig:attribute-binding} illustrates capture and derivation through
a separate transpose-absorption example. The input graph supplies \code{@tr}
and \code{@d}; a user-written \code{f_transform} computes \code{@new_d}.
The table distinguishes the two captured descriptors from the derived one.
Its short fields \code{perm}, \code{l_dim}, and \code{r_dim} belong to a
schematic dialect, rather than the generalized-dot schema used later.
For the displayed permutation \code{[1,0,2]}, the left contraction position
changes from \code{1} in the transposed operand to \code{0} in the original
operand. The figure illustrates descriptor dataflow for this instance; an
arbitrary transpose-absorption rule still requires semantic checking.
\end{figuretext}

\FloatBarrier
\subsection{Writing Attribute Changes}
Conditions and replacement attributes are metadata computations. Expressions
can read captured attributes and available shape and dtype bindings, but they
do not execute tensor computations. Local bindings, builtins, lists, indexing,
conditionals, and comprehensions provide the computations needed to build
descriptors. For transpose composition,
\code{gather(@inner.permutation, @outer.permutation)} computes the replacement
permutation by selecting entries of the inner permutation in outer-permutation
order. Descriptor constructors collect computed fields into immutable
attributes for an RHS operation. Runtime expressibility and semantic
checkability are distinct: some expressions can execute even when the
compiler cannot establish that they preserve a rule's tensor result.

\FloatBarrier
\subsection{Conditions with \texorpdfstring{\code{where}}{where}}
\label{sec:where}
A \code{where} block ends with the condition deciding whether a rule applies.
Any local calculations precede that condition. Only \code{let} bindings end
with a semicolon. Conditions can use captured metadata, builtins, and local
expressions, combining predicates with \code{&&} and \code{||}. For the running
transpose pattern, an illustrative clause is:
\begin{codeblock}[H]
\begin{lstlisting}[language=TEPL]
where {
  let rank = len(@inner.permutation);
  len(@outer.permutation) == rank
}
\end{lstlisting}
\end{codeblock}
The length test illustrates field access and local computation. For a legal
composition of transposes, the equality already follows from operation
legality. A rule can omit \code{where} when no additional condition is needed.

Conditions may express semantic applicability or application policy; TEPL does
not infer profitability decisions. They can also state axis-routing
relationships, for example:
\begin{codeblock}[H]
\begin{lstlisting}[language=TEPL]
@outer.lhs_batching_dimensions == @inner->B
\end{lstlisting}
\end{codeblock}
Its mapping interpretation is introduced in Section~\ref{sec:partitioning}. For a verified semantic rule,
every correctness obligation must follow from operation legality and supported
user-written conditions.

\FloatBarrier
\subsection{Attribute Construction with \texorpdfstring{\code{derive}}{derive}}
The \code{derive} block names replacement descriptors computed from captured
information. Transpose composition combines two captures and an explicit
descriptor constructor:
\begin{codeblock}[H]
\begin{lstlisting}[language=TEPL]
rule compose_transposes_explicit {
  (transpose[@outer] (transpose[@inner] X))
  => (transpose[@combined] X)
  derive {
    @combined = {
      permutation: gather(@inner.permutation, @outer.permutation)
    };
  }
}
\end{lstlisting}
\end{codeblock}
If the inner permutation is $p$ and the outer permutation is $q$, the resulting
permutation is $r[i]=p[q[i]]$. The three parts of a rule have distinct roles:
\code{where} decides applicability, \code{derive} constructs attributes, and
the RHS specifies the replacement computation.

Local bindings can name intermediate metadata computations. A spread such as
\code{{ ..@outer, permutation: permutation }} copies fields before applying
explicit overrides. The \code{where} and \code{derive} blocks have separate
local scopes; derivations can read captured descriptors and earlier derived
descriptors. With the transpose semantics, this explicit constructor is within
the supported semantic-verification fragment. Checking its attribute types
alone would not establish that the rewrite preserves the result.

\FloatBarrier
\subsection{Automatic Attributes with \texorpdfstring{\code{auto}}{auto}}
\label{sec:auto}
Users can omit replacement attributes when operation semantics support
synthesis. An explicit \code{[auto]} annotation requests the same inference:
\begin{codeblock}[H]
\begin{lstlisting}[language=TEPL]
rule compose_transposes_auto {
  (transpose (transpose X)) => (transpose X)
}
\end{lstlisting}
\end{codeblock}
\begin{figuretext}{figures/Introducing auto.pdf}{\code{auto} becomes a reusable attribute function at compilation and is evaluated on matched metadata at runtime.}{fig:introducing-auto}{A vertical flow separates the user rule and operation semantics, compiler
synthesis and verification, and runtime application of the generated function.}
Matching supplies tensor captures and LHS descriptors. Semantic compilation
then derives a reusable program for the RHS permutation. The explicit rule
specifies how to combine two permutations; the automatic rule asks the compiler
to derive that computation from transpose semantics, again obtaining
$r[i]=p[q[i]]$.

Each omitted RHS descriptor is inferred independently. A named annotation
\code{[@name]} instead reuses an existing descriptor, while
\code{[auto, field = value]} pins one field during inference. Users still
provide operation semantics and any required correctness conditions.
Unsupported inference is rejected with a diagnostic. Constructor synthesis
occurs when the rule is compiled; application evaluates the completed
constructor on the matched metadata. Completed constructors can execute in
the greedy rewrite engine or in generated Rust with the shared runtime.
C++ generation remains unavailable.

Figure~\ref{fig:introducing-auto} makes the compile-time/runtime boundary
explicit using matrix transpose absorption. The compiler synthesizes and
checks a function $f(p,q)$ from the rule and operation semantics. Each runtime
match supplies the transpose permutation $p$ and dot descriptor $q$, and the
function supplies the replacement dot descriptor. In the matrix instance,
absorbing permutation \code{[1,0]} changes the left contraction position from
$0$ to $1$. The figure shows only that changing field; it does not imply that
other descriptor fields can be ignored in a complete operation definition.
\end{figuretext}

\FloatBarrier
\subsection{From \texorpdfstring{\code{auto}}{auto} to Attribute Synthesis}
The \code{auto} annotation lets users leave RHS attributes to the compiler,
but the compiler must still determine how to construct them. Matching exposes
known attributes and tensor metadata from the LHS. The challenge is to turn
that information into a program that computes the missing RHS attributes
across ranks and shapes and preserves the computation under the user's
conditions. The next section defines this problem and its correctness contract.

\FloatBarrier
\FloatBarrier
\section{Problem Definition}
\label{sec:problem}
The compiler must construct reusable RHS attribute programs and establish that
the completed rewrite is correct. Table~\ref{tab:notation} gives notation
shared by these requirements. Element types are fixed in this formulation.
\begin{table}[htbp]
\caption{Common notation for attribute synthesis and checking.}
\label{tab:notation}
\begin{tabularx}{\linewidth}{@{}lY@{}}\toprule
Symbol & Meaning \\\midrule
$s$ & Source information: shapes and operation attributes available to the attribute program.\\
$X$ & Input tensor values compatible with the shapes in $s$.\\
$a$ & RHS attributes to construct.\\
$\operatorname{LHS}(s,X)$ & Tensor computed by the LHS.\\
$\operatorname{RHS}(s,a,X)$ & Tensor computed by the RHS using attributes $a$.\\
$\operatorname{Guard}(s)$ & The user-written condition, evaluated successfully to \code{true}.\\\bottomrule
\end{tabularx}
\end{table}

\FloatBarrier
\subsection{What Must the Compiler Find?}
\label{sec:find}
Some RHS attributes are omitted or marked \code{auto}. Finding individual
attribute values at matching time would have the advantage of known ranks,
shapes, and source attributes. TEPL instead derives a formula before the
matched tensors are known. Its inputs are the user-defined operation semantics,
the rewrite rule, and any required conditions. Its output is a reusable program,
rather than an attribute value for one concrete tensor.

Formally, given the operation semantics, rewrite, and condition
$\operatorname{Guard}$, the compiler seeks an attribute program $f$ such that
\begin{equation}
 a=f(s).
 \label{eq:attribute-program}
\end{equation}
The program depends on source information, rather than tensor element values,
and must satisfy Section~\ref{sec:contract}'s contract for every legal source
instance satisfying the guard. For example, composing transposes with
permutations \code{inner} and \code{outer} gives the reusable formula
\code{result[i] = inner[outer[i]]}. Compilation derives the formula without
concrete permutations; matching supplies their values later. An explicit
constructor supplies $f$, while inference constructs it for omitted attributes.

\FloatBarrier
\subsection{What Makes the Result Correct?}
\label{sec:contract}
A completed replacement must have computable attributes, satisfy the RHS
operations' requirements, and produce the same result whenever the rule's
conditions hold. Both explicit constructors and inferred constructors require
checking for semantic rules. The required contract is
\begin{equation}
\begin{aligned}
\forall s.\quad &\operatorname{LegalLHS}(s)\land\operatorname{Guard}(s)\\
\Longrightarrow\quad
 &s\in\operatorname{dom}(f)\ \land\ \operatorname{LegalRHS}(s,f(s))\\
 &{}\land\ \forall X\text{ compatible with }s.\quad
 \operatorname{LHS}(s,X)=\operatorname{RHS}(s,f(s),X).
\end{aligned}
\label{eq:contract}
\end{equation}
In words, every legal LHS satisfying the conditions must yield a computable,
legal RHS with the same result for every compatible input value. Membership
in $\operatorname{dom}(f)$ requires, for example, that inverse permutations
and position lookups be evaluable. Equality includes ordered output shape
and element values under the selected arithmetic policy. Exact algebra alone
does not establish equality under floating-point reassociation.

The guard denotes successful evaluation to \code{true}; it is not an assumption
that a failing computation nevertheless returns a Boolean. The formal account
must respect ordered local evaluation, short-circuiting, indexing, checked
arithmetic, and evaluation failures. These details matter because metadata
programs can be partial. The compiler must establish the contract throughout
the legal guarded domain; otherwise it rejects the semantic rule as unproved.

\FloatBarrier
\subsection{When Can the Compiler Accept a Rewrite?}
The compiler accepts a semantic rewrite when the supported reasoning can show
that replacement attributes are computable, RHS operations are valid, and the
result is preserved under the stated conditions. The same standard applies
to explicit \code{derive} programs and inferred attributes. A rule rejected
as unproved may still be mathematically correct. It can be implemented
manually, but its correctness then has to be established separately.
The checks described below are intended to enforce
Equation~\eqref{eq:contract}; this draft does not promote that intended
contract to a proved soundness theorem for the implementation.

\FloatBarrier
\FloatBarrier
\section{Proposed Method}
\label{sec:method}
\FloatBarrier
\subsection{Axes IR}
\label{sec:axes-ir}
Axes IR is TEPL's intermediate representation of tensor computations. It
describes how operations use tensor axes and compute output values.
Translating operation definitions into this shared representation lets the
compiler track axis relationships across operations, infer replacement
attributes, and check whether completed rewrites preserve the result.

\FloatBarrier
\subsubsection{Symbolic Axis Groups}
\label{sec:axis-groups}
A symbolic axis group represents an ordered sequence of axes without fixing
its length. One group can, for example, describe all contraction axes of an
operation. This representation borrows arbitrary-rank grouping from
TensorRight's aggregated axes~\cite{tensorright}. TensorRight represents an
aggregated axis as a set of named axes; TEPL uses an ordered sequence and
tracks position mappings from groups to tensor axes for attribute construction.

\begin{figuretext}{figures/Symbolic axis groups.pdf}{Ordered, arbitrary-length axis groups preserve operand pairings and determine output order.}{fig:axis-groups}{Four arbitrary-length groups for batch, left-free, contraction, and right-free
axes are connected to operand position lists and an output-axis construction.}
TEPL declares a group with \code{bind}, associating it with axis positions in
one or more operands. When two operands bind the same group, their selected
axes are paired in list order. Attributes determine those positions, allowing
the same definition to describe different ranks and layouts. The two selections
must have equal length and equal sizes at corresponding positions. Matching
sizes alone do not establish pairing: the binding specifies which coordinates
correspond.

Figure~\ref{fig:axis-groups} illustrates arbitrary-length groups, one concrete
dot binding, and output construction. In that instance, the contraction axes
appear in different physical orders in the two operands. Selecting \code{[2,3]}
in \code{X} and \code{[5,2]} in \code{Y} nevertheless identifies the same
ordered contraction group. Reducing it leaves the ordered output
$B\catseq I\catseq J$.


We use four related notions. A tensor's ordered axes form an interface
$P=\langle a_0,\ldots,a_{n-1}\rangle$ in dimension order, together with established
mappings from groups into $P$. Axis labels denote logical coordinates, rather
than sizes. A group $G=\langle g_0,\ldots,g_{m-1}\rangle$ is an ordered
sequence whose order can differ from dimension order. A position sequence
$p=\langle p_0,\ldots,p_{m-1}\rangle$, with $0\leq p_j<n$, locates that
group by $g_j=a_{p_j}$. Positions are zero-based. Finally, $\gath(P,p)$
selects axes in the order specified by $p$, so a mapping records
\begin{equation}
 \gath(P,p)=G.
 \label{eq:mapping}
\end{equation}
\end{figuretext}

The following \code{dot_general} definition declares four attributes and their
axis bindings. Dtype and value semantics are omitted from this excerpt:
\clearpage
\begin{codeblock}[H]
\begin{lstlisting}[language=TEPL]
op dot_general(lhs: tensor, rhs: tensor) -> tensor {
  attrs {
    lhs_batching_dimensions: index[];
    rhs_batching_dimensions: index[];
    lhs_contracting_dimensions: index[];
    rhs_contracting_dimensions: index[];
  }
  axes {
    let lb = attrs.lhs_batching_dimensions;
    let rb = attrs.rhs_batching_dimensions;
    let lc = attrs.lhs_contracting_dimensions;
    let rc = attrs.rhs_contracting_dimensions;
    let lf = ordered_complement(all_axes(lhs), lb ++ lc);
    let rf = ordered_complement(all_axes(rhs), rb ++ rc);
    bind B { lhs: lb, rhs: rb };
    bind K { lhs: lc, rhs: rc };
    bind I { lhs: lf };
    bind J { rhs: rf };
    yield B ++ I ++ J;
  }
}
\end{lstlisting}
\end{codeblock}
Here, \code{all_axes(t)} lists all positions of tensor \code{t}, and \code{++}
concatenates sequences. The builtin \code{ordered_complement} selects remaining
positions in operand order. The axis \code{yield} specifies output order.

As a second concrete instance, smaller than the one in
Figure~\ref{fig:axis-groups}, consider:
\begin{codeblock}[H]
\begin{lstlisting}
lhs:    [i, k0, b, k1]
rhs:    [k1, j, k0, b]
output: [b, i, j]

lhs_batching_dimensions = [2]
rhs_batching_dimensions = [3]
lhs_contracting_dimensions = [1, 3]
rhs_contracting_dimensions = [2, 0]
\end{lstlisting}
\end{codeblock}
The batch binding pairs position $2$ of \code{lhs} with position $3$ of
\code{rhs}. The contraction lists pair $(1,2)$ for $k_0$ and $(3,0)$ for
$k_1$. Remaining axes give $I=\langle i\rangle$ and $J=\langle j\rangle$;
$K$ is reduced and therefore omitted from the output. In mapping notation,
\begin{equation}
\begin{aligned}
P_{\mathrm{lhs}}&=\langle i,k_0,b,k_1\rangle,
 &p_{\mathrm{lhs}}&=\langle1,3\rangle,\\
P_{\mathrm{rhs}}&=\langle k_1,j,k_0,b\rangle,
 &p_{\mathrm{rhs}}&=\langle2,0\rangle,\\
G_{\mathrm{lhs}}&=G_{\mathrm{rhs}}=K=\langle k_0,k_1\rangle.
\end{aligned}
\end{equation}
Consequently,
\begin{equation}
 \gath(P_{\mathrm{lhs}},p_{\mathrm{lhs}})
 =K=\gath(P_{\mathrm{rhs}},p_{\mathrm{rhs}}).
\end{equation}
Although $K$ has two axes here, neither group length nor position values need
be known at rule-compilation time. Grouping preserves the axes separately and
in order; it does not flatten them into one dimension. Groups can be empty
where an operation permits it, such as an unbatched dot's batch group. Their
identities are scoped so that unrelated coordinates with equal extents remain
distinct. Mappings let the compiler recover where each logical group occurs
and use those positions as replacement attributes.

\FloatBarrier
\subsubsection{Expressing Tensor Computations}
\label{sec:value-expressions}
The \code{axes} block defines relationships and output order, while
\code{semantics} defines each output element's value. Axes IR distinguishes
\emph{value expressions}, which compute scalar tensor elements, from
\emph{index expressions}, which calculate the indices used to read them.
Ordinary operand references use coordinates established by bindings. Explicit
index expressions handle calculated accesses, including slices and
convolutions. \code{Read} connects the two kinds of expression by reading an
operand at explicitly calculated indices. A \code{semantics} block returns a
value expression with \code{yield}.

For example, elementwise addition without broadcasting binds every axis of
\code{lhs} to the corresponding axis of \code{rhs}, preserving that group:
\begin{codeblock}[H]
\begin{lstlisting}[language=TEPL]
axes {
  bind S {
    lhs: all_axes(lhs),
    rhs: all_axes(rhs)
  };
  yield S;
}
semantics {
  yield Add(lhs, rhs);
}
\end{lstlisting}
\end{codeblock}
The shared binding requires equal rank and matching sizes. At each coordinate
in $S$, the operand references read corresponding elements and \code{Add}
produces their sum. Table~\ref{tab:value-expressions} lists the value primitives.
In this language, \code{Reduce} means summation; it does not imply support
for other reductions such as maximum.

\begin{table}[htbp]
\caption{Value-expression primitives.}
\label{tab:value-expressions}
\begin{tabularx}{\linewidth}{@{}lY@{}}\toprule
Expression & Meaning\\\midrule
\code{input} & Read an operand element using its axis bindings.\\
\code{Const(c)} & Produce a signed integer constant.\\
\code{Add(a, b)} & Add scalar values.\\
\code{Product(a, b)} & Multiply scalar values.\\
\code{Reduce(G, e)} & Sum $e$ over all coordinate combinations of group $G$.\\
\code{Read(input, indices)} & Read an operand at explicitly calculated indices.\\\bottomrule
\end{tabularx}
\end{table}

Index expressions produce an index sequence with one entry for each accessed
axis. Table~\ref{tab:index-expressions} gives their forms. A slice, for example,
uses
\begin{codeblock}[H]
\begin{lstlisting}[language=TEPL]
Read(input,
  IndexOffset(attrs.starts,
    IndexScale(attrs.strides, Coords(O))))
\end{lstlisting}
\end{codeblock}
Here $O$ is the output group. Per axis, the index expression computes
\code{starts + strides * output_indices}. The \code{Read} expression retrieves
the selected value. The index calculation determines the location, and the
enclosing value expression determines the scalar result.

\begin{table}[htbp]
\caption{Index-expression primitives.}
\label{tab:index-expressions}
\begin{tabularx}{\linewidth}{@{}lY@{}}\toprule
Expression & Meaning\\\midrule
\code{Coords(G)} & Current coordinates of $G$ in group order.\\
\code{IndexAdd(a, b)} & Add index sequences elementwise.\\
\code{Concat(a, b)} & Concatenate index sequences.\\
\code{IndexScale(attrs.strides, indices)} & Scale each index by its corresponding entry.\\
\code{IndexOffset(attrs.starts, indices)} & Add the corresponding offset to each index.\\\bottomrule
\end{tabularx}
\end{table}

Value and index expressions share the groups declared in \code{axes}.
Bindings determine participating operand elements; references bound to the
same group use corresponding coordinates. A reduction ranges over its group
while holding the remaining output coordinates fixed. Because a group can
have variable length, one expression can sum over different numbers of
dimensions. The value semantics completing the previous dot definition are:
\begin{codeblock}[H]
\begin{lstlisting}[language=TEPL]
semantics {
  yield Reduce(K, Product(lhs, rhs));
}
\end{lstlisting}
\end{codeblock}
The product reads coordinates determined by $B,I,J,K$, and the reduction sums
over $K$. Output coordinates belong to $B,I,J$, in the declared output order.
The smaller concrete example therefore computes
\begin{equation}
 \operatorname{out}[b,i,j]
 =\sum_{k_0,k_1}\operatorname{lhs}[i,k_0,b,k_1]
                    \operatorname{rhs}[k_1,j,k_0,b].
\end{equation}

These primitives describe elementwise operations, transpose, summation,
and dot products. Slices and supported convolutions calculate input indices;
broadcasting adds output axes while reusing input values along them. Output
shape remains separate from element computation: a tensor filled with zeros,
for example, still needs a specified shape.

\FloatBarrier
\subsubsection{Reasoning About Composed Operations}
\label{sec:composition}
Value semantics provide a common language for tensor rewrites. Different
operations lower to scalar addition, multiplication, and summation, so laws
over these primitives can relate expressions built from different tensor
operations. TEPL normalizes both sides under the arithmetic laws permitted
by the selected policy.

Consider distributing a generalized dot over elementwise addition without
broadcasting, with identical dot descriptors on both sides:
\begin{equation}
 \operatorname{dot\_general}(X+Y,Z)
 \;\longleftrightarrow\;
 \operatorname{dot\_general}(X,Z)+\operatorname{dot\_general}(Y,Z).
\end{equation}
Addition's local value expression is \code{Add(X, Y)} under group $S$.
The dot's local expression is \code{Reduce(K, Product(lhs, Z))} under its own
batch, contraction, and free groups. Substituting the arithmetic expression
for \code{lhs} gives the expected shape of the computation, but does not
establish how positions in $S$ correspond to $K$ or retain their pairing with
\code{Z}'s contraction axes.

Symbolic axis mapping supplies this missing correspondence. The dot's
contraction attribute selects positions in addition's output. Since addition
preserves axis order, those positions identify corresponding axes in both
\code{X} and \code{Y}. The compiler records this relationship, and similarly
connects batch and free axes, even when group lengths are unknown. With fixed
output coordinates and indexed reads abbreviated, the two expressions become
\begin{equation}
 \sum_K(X+Y)Z
 \qquad\text{and}\qquad
 \sum_K XZ+\sum_K YZ.
\end{equation}
In the value language they are
\begin{codeblock}[H]
\begin{lstlisting}[language=TEPL]
Reduce(K, Product(Add(X, Y), Z))

Add(Reduce(K, Product(X, Z)),
    Reduce(K, Product(Y, Z)))
\end{lstlisting}
\end{codeblock}
Here, operand names abbreviate reads under aligned bindings, and $K$ is the
corresponding contraction group on both sides. Distribution of multiplication
and summation over addition brings the expressions to a common form when
the arithmetic policy permits those laws. Section~\ref{sec:axis-mapping}
explains how the required mappings are represented, composed, and refined.

\FloatBarrier
\subsection{Symbolic Axis Mapping}
\label{sec:axis-mapping}
Symbolic axis mapping connects groups defined by different operations. It
records where each group appears in a tensor, allowing composition of value
expressions while preserving axis order and operand pairings.

\FloatBarrier
\subsubsection{Building and Composing Mapping Tables}
\label{sec:mapping-tables}
Each record $(P,G,p)$ satisfies Equation~\eqref{eq:mapping}. We display records
in tables whose columns denote tensors, whose first row gives the current
interface, and whose other rows give group positions. A dash means that the
example has no mapping for that group. Tables can include both registered
and derivable mappings.

\begin{figuretext}{figures/Symbolic Axis Mapping.pdf}{Shared tensor interfaces connect local axis descriptions without changing tensor values or physical axis order.}{fig:axis-mapping}{Local dot and addition mapping tables are connected through their shared
intermediate, producing a table whose axes use a common coordinate description.}
Figure~\ref{fig:axis-mapping} first illustrates the principle with
\code{T = dot(X,Y);} \code{O = add(T,Z)}. The dot describes its output using $I$ and
$J$, while addition describes that same interface using $S$. Connecting
$S\leftrightarrow I\catseq J$ through $T$ aligns the coordinate descriptions
without changing the arithmetic expression or physical axis order. The
concrete matrix positions below the symbolic expressions are examples,
rather than fixed-rank restrictions on the method. Unlike the following
worked example, the figure places the dot inside the addition; both require
connecting local descriptions through a shared tensor.


We now work through the distribution example
\code{T = Add(X,Y);} \code{O = dot_general(T,Z)} from
Section~\ref{sec:composition}. Let $\rho(G)=\operatorname{range}(|G|)$ denote
all positions in $G$. Ranks and position attributes remain symbolic.
\paragraph{1. Build addition's interface.}
Addition binds \code{X} and \code{Y} to one ordered group $S$ and yields $S$.
Its reads and result thus have these mappings:
\begin{center}
\begin{tabular}{@{}lccc@{}}\toprule
Interface/group & $X$ & $Y$ & $T$\\\midrule
Axis interface & $S$ & $S$ & $S$\\
$S$ & $\rho(S)$ & $\rho(S)$ & $\rho(S)$\\\bottomrule
\end{tabular}
\end{center}

\end{figuretext}

\paragraph{2. Establish the dot's operand layouts.}
The dot introduces batch group $B$, contraction group $K$, and free groups
$I,J$. Let $b_L,b_R$ be its batch attributes, $c_L,c_R$ its contraction
attributes, and $R$ the initial interface of \code{Z}. Free positions are
\begin{equation}
\begin{aligned}
 f_L&=\ocomp(\rho(S),b_L\catseq c_L),\\
 f_R&=\ocomp(\rho(R),b_R\catseq c_R).
\end{aligned}
\end{equation}
Bindings require valid position lists and matching paired extents and must
cover each operand without overlap. Groups can be interleaved, so concatenating
$B,K,I$ does not necessarily reproduce the left operand's dimension order.
Instead, binding positions restore the original order:
\begin{equation}
\begin{aligned}
 \widehat S&=\gath(B\catseq K\catseq I,
                  \operatorname{inverse}(b_L\catseq c_L\catseq f_L)),\\
 \widehat R&=\gath(B\catseq K\catseq J,
                  \operatorname{inverse}(b_R\catseq c_R\catseq f_R)).
\end{aligned}
\label{eq:rebuilt-interfaces}
\end{equation}
Each concatenated position sequence is a permutation because the bindings
form a complete disjoint partition. Thus
\begin{equation}
\begin{aligned}
 \gath(\widehat S,b_L)&=B,&
 \gath(\widehat S,c_L)&=K,&
 \gath(\widehat S,f_L)&=I.
\end{aligned}
\end{equation}
The right interface satisfies analogous bindings. Registering these mappings
and partitions gives
\begin{center}
\begin{tabular}{@{}lcc@{}}\toprule
Interface/group & $T$ & $Z$\\\midrule
Axis interface & $\widehat S$ & $\widehat R$\\
$B$ & $b_L$ & $b_R$\\
$K$ & $c_L$ & $c_R$\\
$I$ & $f_L$ & --\\
$J$ & -- & $f_R$\\\bottomrule
\end{tabular}
\end{center}
The interfaces describe the operands in dot coordinates; they do not change
physical axis order. An existing producer interface is reused when it already
satisfies the bindings.

\paragraph{3. Apply the dot's mappings to \code{X} and \code{Y}.}
Addition describes \code{T} using $S$, whereas the dot describes those same
axes using $\widehat S$. Since addition preserves order, each axis of its
output corresponds to the same-position axes in \code{X} and \code{Y}.
Transporting its reads into $\widehat S$ therefore places $B,K,I$ at positions
$b_L,c_L,f_L$ in all three tensors. The compiler can now compose
\code{Reduce(K, Product(Add(X,Y), Z))}, with $K$ referring to corresponding
coordinates in \code{X}, \code{Y}, and \code{Z}.

\paragraph{4. Extend the table to the output.}
A later consumer sees the dot's output, where surviving axes occur in output
order and $K$ is absent. Registering the order $B\catseq I\catseq J$ extends
the table as follows. The row $\widehat S$ denotes addition's whole-axis group
after coordinate transport.
\begin{center}
\small
\setlength{\tabcolsep}{4pt}
\begin{tabular}{@{}lccccc@{}}\toprule
Interface/group & $X$ & $Y$ & $T$ & $Z$ & $O$\\\midrule
Axis interface & $\widehat S$ & $\widehat S$ & $\widehat S$ & $\widehat R$ & $B\catseq I\catseq J$\\
$\widehat S$ & $\rho(\widehat S)$ & $\rho(\widehat S)$ & $\rho(\widehat S)$ & -- & --\\
$B$ & $b_L$ & $b_L$ & $b_L$ & $b_R$ & $\rho(B)$\\
$K$ & $c_L$ & $c_L$ & $c_L$ & $c_R$ & --\\
$I$ & $f_L$ & $f_L$ & $f_L$ & -- & $|B|+\rho(I)$\\
$J$ & -- & -- & -- & $f_R$ & $|B|+|I|+\rho(J)$\\\bottomrule
\end{tabular}
\end{center}
Adding a length to a position sequence shifts each entry. The entries for
\code{X} and \code{Y} follow from transporting the dot partition into
addition's reads; output entries follow from the declared order.

\paragraph{5. Find positions through existing mappings.}
Not every relationship needs to be stored directly. Suppose $p_X$ locates
$\widehat S$ in \code{X}, and $c_L$ locates $K$ in $\widehat S$:
\begin{equation}
 \gath(P_X,p_X)=\widehat S,\qquad \gath(\widehat S,c_L)=K.
\end{equation}
Composing the selections yields
\begin{equation}
\begin{aligned}
 K&=\gath(\widehat S,c_L)\\
  &=\gath(\gath(P_X,p_X),c_L)\\
  &=\gath(P_X,\gath(p_X,c_L)).
\end{aligned}
\label{eq:gather-composition}
\end{equation}
Thus the $K,X$ entry is $\gath(p_X,c_L)$. For addition, $p_X$ lists all
positions in order, simplifying this expression to $c_L$. Repeated composition
recovers positions through longer chains. Addition's single group makes the
correspondence straightforward; selecting across several producer groups
requires partitioning them to expose finer relationships.

\FloatBarrier
\subsubsection{Symbolic Axis Partitioning}
\label{sec:partitioning}
Figure~\ref{fig:axis-partitioning} illustrates the information retained when a
consumer selects across producer groups. Three positions must be distinguished:
a selected axis's position within its producer group, its slot in the
consumer's contraction list, and its partner's position in the other operand.

\paperfigure{1}{figures/Symbolic Axis Partitioning (Dot Reassociation).pdf}{Partitioning selected axes by producer origin preserves their contraction partners for reassociation.}{fig:axis-partitioning}{Three panels show selecting producer output axes, splitting them by origin,
and retaining contraction partners, followed by a subgroup mapping table.}

\paragraph{1. Why reassociation requires splitting groups.}
Dot reassociation changes $(XY)Z$ to $X(YZ)$. For generalized dots it must also
specify the axes contracted by each replacement. Consider this rule, using
\code{dot} for a dialect's generalized dot:
\begin{codeblock}[H]
\begin{lstlisting}[language=TEPL]
rule associate {
  (dot[@outer] (dot[@inner] X Y) Z)
  => (dot[auto] X (dot[auto] Y Z))
  where {
    @outer.lhs_batching_dimensions == @inner->B
  }
}
\end{lstlisting}
\end{codeblock}
Incompatible batch routing can change a reassociation's result. The guard
requires the outer dot's left batch list to equal the positions of the inner
dot's batch group in its result, in the same order. The expression
\code{@inner->B} denotes those positions. This restricts applicability;
the completed replacement still requires the checks in
Sections~\ref{sec:synthesis} and~\ref{sec:correctness}.

Let $T=\operatorname{dot\_general}(X,Y)$. The inner dot contracts $K$ and
yields $P_T=B\catseq I\catseq J$, with free groups $I$ from \code{X} and
$J$ from \code{Y}. Write $c_L^{\mathrm{outer}},c_R^{\mathrm{outer}}$ for the
outer contraction lists in \code{T} and \code{Z}. Their shared group is
\begin{equation}
 L=\gath(P_T,c_L^{\mathrm{outer}}).
\end{equation}
Since $B$ remains batch, the left contraction selects only from $I,J$.
Unlike addition's single group, this producer has groups from different
operands. The outer dot can select part of each, mixing the selected axes
in $L$. To build \code{dot_general(Y,Z)}, the compiler must separate the
portion from $J$, which can be contracted there, from the partners of $I$,
which must survive for the final dot with \code{X}.

\paragraph{2. Find the selected part of each producer group.}
For each producer group, retain both within-group positions and contraction
slots. Group $I$ occupies
\begin{equation}
 p_I=|B|+\rho(I)
\end{equation}
in $P_T$. Compare this list to the outer selection using
\begin{equation}
 (u_I,v_I)=\mpo(c_L^{\mathrm{outer}},p_I).
\end{equation}
For distinct-entry sequences $a,b$, the operation
$\mpo(a,b)=(u,v)$ scans $a$ in order. Whenever $a_i=b_j$, it appends $i$ to
$u$ and $j$ to $v$. It returns positions in the input lists, rather than
the shared values. Hence $\gath(a,u)=\gath(b,v)$, with both results empty
when there is no overlap. For example,
\begin{equation}
\begin{aligned}
 a&=\langle3,0,4\rangle,&b&=\langle0,2,3\rangle,\\
 u&=\langle0,1\rangle,&v&=\langle2,0\rangle.
\end{aligned}
\end{equation}
The two gathers produce $\langle3,0\rangle$, retaining $a$'s order.
For symbolic inputs, the results remain position expressions until they can
be evaluated or simplified.

Here $u_I$ locates the selected axes in the contraction list, and $v_I$
locates them within $I$:
\begin{equation}
 \gath(c_L^{\mathrm{outer}},u_I)=\gath(p_I,v_I).
\end{equation}
The selected subgroup is therefore
\begin{equation}
 \Isel=\gath(I,v_I)=\gath(L,u_I).
\end{equation}
The remaining view uses an ordered complement:
\begin{equation}
 w_I=\ocomp(\rho(I),v_I),\qquad \Irem=\gath(I,w_I).
\end{equation}
The selected view follows outer contraction order, whereas the remaining
view follows $I$'s original order. Their concatenation need not reproduce
$I$, so the position sequences must be retained. The same construction
applies to $J$ using $p_J=|B|+|I|+\rho(J)$. For either $G\in\{I,J\}$,
$u_G,v_G$ locate its selected portion, and
$w_G=\ocomp(\rho(G),v_G)$ locates its remaining portion. The lengths and
selected portions remain symbolic.

Partitioning adds the following relationships:
\begin{center}
\small
\begin{tabular}{@{}lcccc@{}}\toprule
Group & In $P_T$ & In $I$ & In $J$ & In $L$\\\midrule
$\Isel$ & $\gath(p_I,v_I)$ & $v_I$ & -- & $u_I$\\
$\Irem$ & $\gath(p_I,w_I)$ & $w_I$ & -- & --\\
$\Jsel$ & $\gath(p_J,v_J)$ & -- & $v_J$ & $u_J$\\
$\Jrem$ & $\gath(p_J,w_J)$ & -- & $w_J$ & --\\\bottomrule
\end{tabular}
\end{center}
The $L$ column supplies the connection to \code{Z} through contraction slots.

\paragraph{3. Map selected axes to \code{Z}.}
The outer dot pairs its left and right lists entry by entry. Selecting slots
$u_I$ and $u_J$ from its right list therefore gives the positions in \code{Z}
paired with the selected subgroups:
\begin{equation}
 z_I=\gath(c_R^{\mathrm{outer}},u_I),\qquad
 z_J=\gath(c_R^{\mathrm{outer}},u_J).
\label{eq:z-partners}
\end{equation}
In Figure~\ref{fig:axis-partitioning}, \code{Z} has axes $[z_0,f,z_1]$ and
right contraction positions $[2,0]$. The left list selects $[j_1,i_0]$.
Axis $j_1$ occupies contraction slot $0$ and pairs with $z_1$ at position $2$;
$i_0$ occupies slot $1$ and pairs with $z_0$ at position $0$.
The selected groups thus retain all three descriptions:
\begin{center}
\begin{tabular}{@{}lccc@{}}\toprule
Selected group & In producer group & Slots in $L$ & In $Z$\\\midrule
$\Isel$ & $v_I$ & $u_I$ & $z_I$\\
$\Jsel$ & $v_J$ & $u_J$ & $z_J$\\\bottomrule
\end{tabular}
\end{center}
Combining these entries with producer mappings describes the original
computation in terms of the refined groups without losing operand pairings.

\paragraph{4. Extend the original mapping table.}
The original graph is $T=\operatorname{dot\_general}(X,Y)$ followed by
$O=\operatorname{dot\_general}(T,Z)$. Write $b_X,k_X,i_X$ for the known
positions of $B,K,I$ in \code{X}, and $b_Y,k_Y,j_Y$ for those of $B,K,J$ in
\code{Y}. Composition locates each subgroup: for example,
$\gath(i_X,v_I)$ locates $\Isel$ in \code{X}. In \code{Z}, let $b_Z$ be the
batch positions and $f_Z$ the positions of free group $F$.

The inner reduction removes $K$ from $T$; the outer reduction removes
$\Isel,\Jsel$ from $O$. The surviving order is
\begin{equation}
 P_O=B\catseq\Irem\catseq\Jrem\catseq F.
\label{eq:source-output}
\end{equation}
Table~\ref{tab:refined-mappings} combines the refined relationships. To keep
entries readable, define
\begin{equation}
\begin{aligned}
 i_s&=\gath(i_X,v_I),&i_r&=\gath(i_X,w_I),\\
 j_s&=\gath(j_Y,v_J),&j_r&=\gath(j_Y,w_J),\\
 t_{Is}&=\gath(p_I,v_I),&t_{Ir}&=\gath(p_I,w_I),\\
 t_{Js}&=\gath(p_J,v_J),&t_{Jr}&=\gath(p_J,w_J),\\
 o_I&=|B|+\rho(\Irem),\\
 o_J&=|B|+|\Irem|+\rho(\Jrem),\\
 o_F&=|B|+|\Irem|+|\Jrem|+\rho(F).
\end{aligned}
\end{equation}
These are abbreviations for the position expressions, not additional
attributes to be supplied by the user.
\begin{table}[htbp]
\caption{Refined mappings across the original reassociation graph. Each entry
gives positions in the column's interface. A dash indicates no mapping in
this example.}
\label{tab:refined-mappings}
\centering
\begin{tabular}{@{}lccccc@{}}\toprule
Group/interface & $X$ & $Y$ & $T$ & $Z$ & $O$\\\midrule
Axis interface & $P_X$ & $P_Y$ & $P_T$ & $P_Z$ & $P_O$\\
$B$ & $b_X$ & $b_Y$ & $\rho(B)$ & $b_Z$ & $\rho(B)$\\
$K$ & $k_X$ & $k_Y$ & -- & -- & --\\
$\Isel$ & $i_s$ & -- & $t_{Is}$ & $z_I$ & --\\
$\Irem$ & $i_r$ & -- & $t_{Ir}$ & -- & $o_I$\\
$\Jsel$ & -- & $j_s$ & $t_{Js}$ & $z_J$ & --\\
$\Jrem$ & -- & $j_r$ & $t_{Jr}$ & -- & $o_J$\\
$F$ & -- & -- & -- & $f_Z$ & $o_F$\\\bottomrule
\end{tabular}
\end{table}

The procedure generalizes to other producer--consumer pairs: compare position
sequences on the shared interface, identify overlap, and add mappings from
each overlapping portion into both parent groups. Each partition is internally
disjoint, while groups in different partitions can overlap.

\FloatBarrier
\subsubsection{General Construction Algorithm}
\label{sec:mapping-algorithm}
The construction algorithm processes producers before consumers, connects each
operation's bindings to its inputs, and extends mappings to the output.
Overlap introduces the finer subgroup relationships described above.
The input is a source graph and its operation definitions; the output comprises
tensor interfaces, composed value expressions, established mappings,
and partitions.

\paragraph{State.}
The state is $\Sigma=(V,M,\Pi)$. For each tensor $t$,
$V[t]=(P_t,e_t)$ stores its ordered interface and value expression.
$M$ contains individual mapping entries $(P,G,p)$. The partition collection
$\Pi$ records which groups cover an interface completely and without overlap:
\begin{equation}
 \pi=\bigl(P,\langle(G_1,p_1),\ldots,(G_m,p_m)\rangle\bigr).
\end{equation}
An interface can have several partitions because different operations group
its axes differently. A partition changes the grouping used in reasoning;
it does not change tensor shape, rank, or axis order. $M$ locates the groups,
whereas $\Pi$ says which groups jointly cover the interface exactly once.

For example, if addition produces $T$ with $S=\langle i,j\rangle$, then
$V[T]=(P_T,\operatorname{Add}(X,Y))$ with $P_T=S$,
$M$ includes $(P_T,S,\langle0,1\rangle)$, and $\Pi$ includes the single-member
partition $(P_T,\langle(S,\langle0,1\rangle)\rangle)$. A consumer may split
$S$ into $I=\langle i\rangle$ and $J=\langle j\rangle$. Then $M$ gains their
position mappings and $\Pi$ gains the two-member partition alongside the
original one. A shape $\langle d_i,d_j\rangle$ remains unchanged.

\paragraph{Invariants.}
Every mapping must select its group in the stated order:
\begin{equation}
 (P,G,p)\in M\ \Longrightarrow\ \gath(P,p)=G.
 \label{eq:mapping-invariant}
\end{equation}
Every partition member satisfies the same invariant, and
$p_1\catseq\cdots\catseq p_m$ must be a permutation of $\rho(P)$.
Thus each axis occurs exactly once within one partition, although different
partitions may overlap. In the addition example,
$\gath(\langle i,j\rangle,\langle0,1\rangle)=S$, and $\langle0,1\rangle$
covers every position. The sequence $\langle0,0\rangle$ would repeat one
axis and omit the other, violating the partition invariant.

\begin{paperalgorithm}[H]
\caption{Build symbolic axis mapping.}\label{alg:mapping}
\begin{lstlisting}[language=Python]
def build_symbolic_axis_mapping(graph):
    V = initialize_inputs(graph)
    M = set()       # Entries (interface, group, positions)
    partitions = {}

    for op in producer_first(graph):
        inputs = [V[t] for t in op.inputs]
        groups, bindings = resolve_bindings(op, inputs)
        aligned = []

        for (P, expression), operand_bindings in zip(inputs, bindings):
            Q = build_interface(P, operand_bindings)
            for group, positions in operand_bindings:
                M.add((Q, group, positions))
            refine_and_register(Q, operand_bindings, M, partitions)
            expression = transport(expression, P, Q, M, partitions)
            aligned.append((Q, expression))

        Pout = evaluate(op.output_axes, groups)
        output_bindings = list(output_positions(op, groups))
        for group, positions in output_bindings:
            M.add((Pout, group, positions))
        refine_and_register(Pout, output_bindings, M, partitions)
        V[op] = (Pout, compose(op.value, aligned, groups))

    return V, M, partitions
\end{lstlisting}
\end{paperalgorithm}

\code{initialize_inputs} gives each graph input an interface and read
expression. \code{producer_first} ensures descriptions are available before
use. \code{resolve_bindings} evaluates symbolic position expressions and
establishes validity and pairing conditions. It reuses producer coordinates
where applicable and otherwise introduces local groups. Each operand's
bindings $(G_i,p_i)$ cover its axes.

\code{build_interface} reuses $P$ when the required groups already occupy the
required positions. Otherwise it constructs
\begin{equation}
 Q=\gath\bigl(G_1\catseq\cdots\catseq G_m,
        \operatorname{inverse}(p_1\catseq\cdots\catseq p_m)\bigr).
\end{equation}
The inverse arranges collected groups in tensor dimension order.
\code{refine_and_register} registers the new partition and compares it
with earlier partitions on the same interface. It uses \code{matchPositions}
to add overlap mappings. \code{transport} updates the producer's read
coordinates from $P$ to $Q$ and carries mappings and partitions into that
description for the current consumer. Equal interfaces need no coordinate
update. Finally, \code{evaluate} and \code{output_positions} derive the
output interface and positions from the declared order, and \code{compose}
combines transported expressions using the operation's value definition.

\paragraph{Deriving mappings on demand.}
A missing mapping can be recovered by following established relationships:
\begin{equation}
\begin{aligned}
 (P,G,p)\in M\ \land\ (G,H,q)\in M
 \quad\Longrightarrow\quad
 \gath(P,\gath(p,q))=H.
\end{aligned}
\end{equation}
Lookup can return $\gath(p,q)$ or register the corresponding entry. Some
relationships propagate during construction and others are derived when
queried; enumerating every possible mapping in advance is unnecessary.
Local bindings and output order establish local mappings, paired position
matches establish overlaps, and Equation~\eqref{eq:gather-composition}
composes them. These operations are designed to maintain the invariant;
they do not dispense with checking the completed rewrite.

\FloatBarrier
\subsection{Attribute Program Synthesis}
\label{sec:synthesis}
The mapping procedure establishes which logical axes occur in source tensors
and where they occur. Synthesis uses these relationships to fill the
replacement's attribute holes. It normalizes value expressions, identifies
which axes each replacement operation must preserve or reduce, and queries
mappings to construct its attributes.

Figure~\ref{fig:auto-matching} previews this process for distribution of a dot
over addition. The source is lowered into Axes IR and normalized; the RHS
arithmetic skeleton supplies reads to match before its attributes are known.
Read correspondence allows reuse of source coordinates, after which axis
positions construct the two missing dot descriptors. Here addition has no
broadcasting and preserves paired axis order, so both inferred descriptors
equal the source descriptor. The figure's reads $A[I,K]$, $B[I,K]$, and
$C[K,J]$ abbreviate group coordinates, rather than asserting that groups are
contiguous in physical layouts. The skeleton's question marks denote unresolved
reduction information. Equivalence is checked only after construction.

\paperfigure{1}{figures/Auto attribute matching_ from rules to attributes.pdf}{Matching normalized operand reads supplies coordinate correspondences for constructing replacement attributes.}{fig:auto-matching}{Three stages connect a distribution rule and operation semantics to normalized
terms, read matching, position queries, and completed replacement descriptors.}

For attribute reconstruction beyond descriptor reuse, continue the guarded
reassociation example under an arithmetic policy permitting the required laws:
\begin{equation}
 \underbrace{\operatorname{dot}(\operatorname{dot}(X,Y),Z)}_{\mathrm{LHS}}
 \ \longrightarrow\
 \underbrace{\operatorname{dot}(X,\operatorname{dot}(Y,Z))}_{\mathrm{RHS}}.
\end{equation}
The rule supplies the replacement graph. Let $T=\operatorname{dot}(X,Y)$ be
the source intermediate, $U=\operatorname{dot}(Y,Z)$ the replacement
intermediate, and $O,O'$ their respective final results.

\FloatBarrier
\subsubsection{Value Expression Normalization}
\label{sec:normalization}
The graphs combine their reads in different orders. Normalization identifies
corresponding computations before constructing the replacement's attributes.
With output coordinates fixed and indexed reads abbreviated, the LHS is
\begin{equation}
 e_{\mathrm{LHS}}=\sum_L\left(\sum_K XY\right)Z.
\end{equation}
Each sum ranges over the coordinates of its group. The abbreviated notation
hides, but does not discard, the mappings established in
Section~\ref{sec:axis-mapping}.

TEPL uses a sum-of-products representation similar to SPORES's R-polynomial
form~\cite[Section~2.3, Definition~2.4]{spores}. A term contains a coefficient,
indexed operand reads, and symbolic reduction axes:
\begin{equation}
 t=(c,[r_1,\ldots,r_m],R),\qquad
 \operatorname{value}(t)=\sum_R c\prod_{i=1}^{m}r_i.
 \label{eq:normalized-term}
\end{equation}
Each read retains its tensor identity and coordinate expression. $R$ is an
ordered sequence of symbolic axes, and an empty sequence denotes no summation.
The empty product is $1$, permitting constants. A normalized expression is a
list $N=[t_1,\ldots,t_n]$ denoting
\begin{equation}
 \operatorname{value}(N)=\sum_{j=1}^{n}\operatorname{value}(t_j).
\end{equation}
Reads within a term multiply, separate terms add, and the empty term list
represents zero. Examples, again abbreviating indexed reads, are:
\begin{center}
\begin{tabular}{@{}ll@{}}\toprule
Expression & Normalized term list\\\midrule
$XY$ & \code{[(1, [X, Y], [])]}\\
$X(Y+Z)=XY+XZ$ & \code{[(1, [X, Y], []), (1, [X, Z], [])]}\\
$2XY-Z$ & \code{[(2, [X, Y], []), (-1, [Z], [])]}\\
$\sum_K XY$ & \code{[(1, [X, Y], K)]}\\\bottomrule
\end{tabular}
\end{center}
This similarity to SPORES concerns representation; it does not transfer a
completeness result to TEPL's supported reasoning.

Under the permitted multiplication and reduction laws, the example normalizes
to
\begin{equation}
 \sum_L\left(\sum_K XY\right)Z
 =\sum_L\sum_K XYZ
 =\sum_{L,K}XYZ.
\end{equation}
The inner reduction contributes reads \code{X}, \code{Y}, and axes $K$.
Multiplication incorporates \code{Z} after checking that it does not depend
on $K$. The outer reduction prepends $L$ after establishing that $L$ and $K$
are disjoint. Thus
\begin{codeblock}[H]
\begin{lstlisting}
Value expression:
Reduce(L,
  Product(
    Reduce(K, Product(X, Y)),
    Z))

Normalized term:
t = (1, [X, Y, Z], L ++ K)
\end{lstlisting}
\end{codeblock}

Unknown RHS attributes mean its actual reduction axes and read coordinates
are not yet available. Synthesis instead normalizes its \emph{arithmetic
skeleton}: it expands declared value operations while temporarily omitting
axis information. The skeleton for $X(YZ)$ has coefficient $1$, reads
\code{[X,Y,Z]}, and an empty reduction sequence. That empty sequence is a
property of the skeleton representation, not a claim that the completed RHS
has no reductions.

Synthesis matches each RHS read with the same operand's LHS read and reuses
its coordinates. For several terms, it first matches terms with the same
operands. Multiple possible read-coordinate assignments must be shown to
agree; otherwise synthesis remains unresolved. This correspondence provides
coordinates for construction. Reduction domains, coefficients, and full
value equivalence still have to be checked after attributes are known.

\FloatBarrier
\subsubsection{Determine Which Axes Each Operation Preserves}
\label{sec:keep-axes}
After read matching, the logical axes in each RHS input are known, but their
roles in the replacement may change. The source inner dot contracts $K$;
the source outer dot contracts $L$, partitioned into $\Isel$ and $\Jsel$.
The final output retains $B,\Irem,\Jrem,F$. Source mappings show that
\code{Y} contains $B,K,\Jsel,\Jrem$ and \code{Z} contains
$B,\Isel,\Jsel,F$. These lists specify membership rather than physical order.
To build $U$, the compiler must classify these axes as batch, contraction,
left-free, or right-free. In particular, $K$ must survive $U$ for a later
contraction with \code{X}.

\begin{figuretext}{figures/Infer Axis Roles in U.pdf}{An intermediate preserves axes needed outside its computation or in the final output and reduces the remaining shared axes.}{fig:infer-u}{A flow compares axes inside the replacement intermediate with axes needed
outside it and in the output, retaining required axes and contracting the rest.}
Figure~\ref{fig:infer-u} shows the decision for an unbatched instance.
Starting with axes in \code{Y} or \code{Z}, it first preserves those needed
by \code{X}, then those required by the final output. The remaining shared
group, $\Jsel$, can be reduced inside $U$. The figure omits batch group $B$;
the equations below retain it for the guarded general case.

\paragraph{General rule for preserving axes.}
An intermediate must keep an axis used to compute it if that axis is also
needed outside the computation or in the final output. For intermediate $t$,
its replacement subtree includes the operation and its input reads. In $U$,
reads of \code{Y} and \code{Z} are inside; the read of \code{X} is outside.
Let $\Axes(x)$ denote the set of logical axes in tensor $x$'s interface,
ignoring order. In a matched product, let $\Axes_{\mathrm{in}}(t)$ and
$\Axes_{\mathrm{out}}(t)$ collect axes used by reads inside and outside the
subtree. With original output $O$, the required set is
\begin{equation}
\begin{aligned}
 \Keep(t)&=\Axes_{\mathrm{in}}(t)\cap\\
 &\quad\bigl(\Axes_{\mathrm{out}}(t)\\
 &\qquad{}\cup\Axes(O)\bigr).
\end{aligned}
 \label{eq:keep-axes}
\end{equation}
These are sets of logical coordinates, with order determined later. For
multiple matched products, the required sets must agree.
Reducing an axis removes its coordinate from the intermediate. An outside
read using that coordinate must still pair its values with the intermediate
at matching coordinates; an output axis must remain available to distinguish
output elements. These requirements give the union in
Equation~\eqref{eq:keep-axes}.
\end{figuretext}

Intersecting with the inside set excludes axes
used only outside $t$, which need not pass through it. Axes used only inside
can be reduced there if the operation's declared semantics permits it.

For $U$, the inside set is $\Axes(Y)\cup\Axes(Z)$ and the outside set is
$\Axes(X)$. Using group names as shorthand for all their constituent axes,
\begin{equation}
\begin{aligned}
 \Axes_{\mathrm{in}}(U)
    &=\{B,K,\Isel,\Jsel,\Jrem,F\},\\
 \Axes_{\mathrm{out}}(U)
    &=\{B,K,\Isel,\Irem\},\\
 \Axes(O)&=\{B,\Irem,\Jrem,F\},\\
 \Keep(U)&=\{B,K,\Isel,\Jrem,F\}.
\end{aligned}
\end{equation}
Group $K$ occurs inside and outside and so is preserved. $\Jsel$ is needed
only inside and can be contracted. Although $\Irem$ belongs to the final
output, it occurs outside $U$ and reaches the final dot through \code{X}.
The resulting decisions are summarized in Table~\ref{tab:keep-decisions}.
At the final dot, $K$ and $\Isel$ are shared by \code{X} and $U$ but absent
from $O$, so they are reduced together.

\begin{table}[htbp]
\caption{Preservation decisions for groups present in $Y$ or $Z$.}
\label{tab:keep-decisions}
\begin{tabularx}{\linewidth}{@{}lYY@{}}\toprule
Group & Need outside $U$ or in $O$ & Action in $U$\\\midrule
$B$ & Used by $X$ and present in $O$ & Preserve as batch axes.\\
$K$ & Used by $X$ & Preserve for final contraction.\\
$\Isel$ & Used by $X$; paired axes occur in $Z$ & Preserve for final contraction.\\
$\Jsel$ & Absent from $X$ and $O$ & Contract between $Y$ and $Z$.\\
$\Jrem,F$ & Present in $O$ & Preserve as free axes.\\\bottomrule
\end{tabularx}
\end{table}

\paragraph{From required axes to operation groups.}
The required set identifies surviving coordinates. Assigning coordinates to
operation groups additionally uses which operands contain them and whether
the declared value expression reduces the group. For an operation producing
$t$ from $x_1,\ldots,x_m$, let $\Axes_i=\Axes(x_i)$, and let $E_g$ be the
nonempty set of operands bound to declared group $g$. Axes appearing in exactly
those operands form
\begin{equation}
 D_g=\left(\bigcap_{i\in E_g}\Axes_i\right)
     \setminus\left(\bigcup_{i\notin E_g}\Axes_i\right).
\end{equation}
The candidate assignment is
\begin{equation}
 \Axes_g=\begin{cases}
 D_g\setminus\Keep(t),&\text{if the value definition reduces }g,\\
 D_g\cap\Keep(t),&\text{otherwise.}
 \end{cases}
 \label{eq:group-assignment}
\end{equation}
This equation assigns membership, with ordered positions constructed afterward.
For a dot with operands $a,b$, shared axes in $\Keep(t)$ are batch, and shared
axes outside $\Keep(t)$ are contracted. Left-only and right-only axes in
$\Keep(t)$ form the respective free groups. Thus the assignments are:
\begin{center}
\small
\begin{tabular}{@{}lcccc@{}}\toprule
Operation & Batch & Contracted & Left free & Right free\\\midrule
$U=\operatorname{dot}(Y,Z)$ & $B$ & $\Jsel$ & $K,\Jrem$ & $\Isel,F$\\
$O'=\operatorname{dot}(X,U)$ & $B$ & $K,\Isel$ & $\Irem$ & $\Jrem,F$\\\bottomrule
\end{tabular}
\end{center}
Free axes follow operand order. The rule relies on distinguishable declared
group roles; ambiguous assignments remain unresolved. Groups introduced by an
explicit output domain rather than operand bindings are handled separately
in Section~\ref{sec:synthesis-algorithm}.

\FloatBarrier
\subsubsection{Construct Attributes from the Mapping Table}
\label{sec:attribute-construction}
Group membership says which axes a dot uses, but its attributes need ordered
positions in actual operands. A query $\pos(P,G)$ returns $p$ satisfying
$\gath(P,p)=G$, using existing mappings and their compositions. Both sides
of a paired binding query the same \emph{ordered} group, ensuring that
corresponding entries select paired axes.

\begin{figuretext}{figures/Constructing Dot Attributes from Axis Mappings.pdf}{Paired position queries construct dot attributes; the completed intermediate supplies mappings for the final dot.}{fig:attribute-construction}{Paired queries on Y and Z construct the inner dot, an interface bridge
establishes U's layout, and paired queries on X and U construct the final dot.}
Figure~\ref{fig:attribute-construction} shows the two stages for the unbatched
example. The first pair of queries constructs the inner replacement's
contraction lists. Its completed interface $P_U$ then supplies mappings for
the final dot's queries. In the figure, $G$ abbreviates
$\Isel\catseq K$, and preserved-group chips specify membership, not order.
The short fields \code{lhs_contract} and \code{rhs_contract} abbreviate the
contraction attributes. The following derivation also includes batch lists.


\paragraph{1. Construct the inner replacement dot.}
For $U=\operatorname{dot}(Y,Z)$, the source table already locates $B$ and
$\Jsel$ in both operands. Its attributes are
\begin{equation}
\begin{aligned}
 b_L^U&=b_Y,&b_R^U&=b_Z,\\
 c_L^U&=\gath(j_Y,v_J),&c_R^U&=z_J.
\end{aligned}
\end{equation}
By Equation~\eqref{eq:z-partners}, $z_J=\gath(c_R^{\mathrm{outer}},u_J)$,
so both lists select $\Jsel$ in the same order. Free positions are
\begin{equation}
\begin{aligned}
 f_Y^U&=\\
 &\ocomp(\rho(P_Y),b_Y\catseq c_L^U),\\
 f_Z^U&=\\
 &\ocomp(\rho(P_Z),b_Z\catseq c_R^U),
\end{aligned}
\end{equation}
and the dot's output order constructs
\begin{equation}
\begin{aligned}
 P_U&=B\catseq\gath(P_Y,f_Y^U)\\
 &\catseq\gath(P_Z,f_Z^U).
\end{aligned}
 \label{eq:u-interface}
\end{equation}
The compiler extends mappings to $P_U$ using this layout. This is necessary
because $K$ and $\Jrem$ can be interleaved in \code{Y}, as can $\Isel$ and
$F$ in \code{Z}. The interface preserves those physical orders instead of
assuming that each surviving group is contiguous.

\end{figuretext}

\paragraph{2. Construct the final replacement dot.}
The final dot contracts $K$ and $\Isel$. Filtering source reduction order
$L\catseq K$ to those axes gives $\Isel\catseq K$: axes from $I$ occur
within $L$, before $K$. Use that same order for both contraction lists:
\begin{equation}
\begin{aligned}
 c_L^{O'}&=\gath(i_X,v_I)\catseq k_X,\\
 c_R^{O'}&=\pos(P_U,\Isel)\catseq\pos(P_U,K).
\end{aligned}
\end{equation}
A fixed binding determines group order, and inferred paired lists must follow
it. Otherwise the implementation prefers source reduction order for reduced
groups and source output order for retained groups, appending remaining axes
in operand order. Table~\ref{tab:replacement-attributes} completes the two
descriptors.
\begin{table}[htbp]
\caption{Candidate attributes for the replacement dots.}
\label{tab:replacement-attributes}
\begin{tabularx}{\linewidth}{@{}lYY@{}}\toprule
Attribute & $U=\operatorname{dot}(Y,Z)$ & $O'=\operatorname{dot}(X,U)$\\\midrule
Left batch & $b_Y$ & $b_X$\\
Right batch & $b_Z$ & $\rho(B)$\\
Left contraction & $\gath(j_Y,v_J)$ & $\gath(i_X,v_I)\catseq k_X$\\
Right contraction & $z_J$ & $\pos(P_U,\Isel)$\newline $\catseq\pos(P_U,K)$\\\bottomrule
\end{tabularx}
\end{table}
Removing the final contraction axes leaves
\begin{equation}
 P_{O'}=B\catseq\Irem\catseq\Jrem\catseq F=P_O.
\end{equation}
These attributes therefore have the candidate pairings and output order.
Operation legality and value equivalence still require checking.

\FloatBarrier
\subsubsection{Produce and Verify Attribute Programs}
\label{sec:emit-programs}
The tables contain symbolic group queries, while a rule application must
compute concrete attributes. The compiler resolves each query to expressions
over source attributes and shapes. For the inner replacement's right
contraction list, for example,
\begin{equation}
\begin{aligned}
 u_J&=\mpo\bigl(c_L^{\mathrm{outer}},
                   |B|+|I|+\rho(J)\bigr)_1,\\
 c_R^U&=\gath(c_R^{\mathrm{outer}},u_J),
\end{aligned}
\label{eq:source-only-example}
\end{equation}
where subscript $1$ selects the first returned sequence. Group lengths and
ranges are themselves derived from source metadata. This expression constructs
positions at each match without accessing tensor elements. Queries involving
$P_U$ similarly substitute its layout and compose mappings until only source
metadata remains. Emitted programs contain neither unresolved RHS attributes
nor logical-group queries.

For calculated accesses, position mappings alone need not determine every
attribute. The implementation additionally matches supported index expressions,
including coordinate coefficients and offsets. Composed slices, for example,
have the candidate attributes
\begin{equation}
\begin{aligned}
 \mathrm{start}'&=\mathrm{start}_{\mathrm{inner}}
   +\mathrm{stride}_{\mathrm{inner}}\,\mathrm{start}_{\mathrm{outer}},\\
 \mathrm{stride}'&=\mathrm{stride}_{\mathrm{inner}}\,
                  \mathrm{stride}_{\mathrm{outer}}.
\end{aligned}
\end{equation}
These expressions must also satisfy resulting access bounds. Coupled or
inverse equations outside the supported fragment remain unresolved.

\paragraph{Check the completed replacement.}
With attributes filled, the RHS has actual coordinates and reduction domains.
In the example, $U$ reduces $\Jsel$ and the final dot reduces
$\Isel\catseq K$:
\begin{equation}
\begin{aligned}
 e_{\mathrm{RHS}}
 &=\sum_{\Isel,K}X\left(\sum_{\Jsel}YZ\right)\\
 &=\sum_{\Isel,K,\Jsel}XYZ.
\end{aligned}
\end{equation}
The selected subgroups partition $L$, and the mappings preserve coordinate
pairing. The normalized RHS can therefore be compared to
$\sum_{L,K}XYZ$ under the permitted reduction laws. Verification also checks
output order, shape, dtype, and operation legality.

After construction, a fresh verification context reconstructs the completed
rule using the emitted programs. Required conditions must follow from legal
source metadata and the declared guard; construction obligations do not
become new assumptions. On success, the result is a collection of checked
source-dependent programs, compiled once and evaluated on each matched
graph's metadata. Unestablished construction or verification relationships
leave the rule unresolved.

\FloatBarrier
\subsubsection{General Synthesis Algorithm}
\label{sec:synthesis-algorithm}
The procedure uses each operation's declared groups, value expression, and
output order, rather than requiring the replacement to contain dots. Its
inputs are a rule's source graph, replacement graph with attribute holes,
operation definitions, and declared conditions. It returns checked attribute
programs or an unresolved result.

The state retains $(V,M,\Pi)$ and adds a table $C$ of proposed attribute
expressions. \todo{Confirm operation-qualified keys for $C$: the pseudocode
uses an attribute name as shorthand, but different replacement operations
can have fields with the same name.} In the pseudocode, \code{.axes} and \code{.value} denote the two
components of $V[t]$. Read correspondence supplies required output sets,
and replacement nodes are processed after their inputs. The construction
invariant requires every completed node to have resolved legal attributes,
an interface and value expression consistent with them, and mappings
satisfying Equation~\eqref{eq:mapping-invariant}. Whole-rule equivalence is
checked after completion.

\begin{paperalgorithm}[H]
\caption{Synthesize attribute programs.}\label{alg:synthesis}
\begin{lstlisting}[language=Python]
def synthesize_attribute_programs(rule):
    V, M, partitions = build_symbolic_axis_mapping(rule.lhs)
    source = normalize(V[rule.lhs].value)
    skeleton = normalize(arithmetic_skeleton(rule.rhs))
    match = correspond(source, skeleton)
    required = required_output_axes(match, V[rule.lhs].axes)
    bind_replacement_reads(V, match)
    C = {}

    for op in producer_first(rule.rhs):
        if op.is_input:
            continue
        inputs = [V[t] for t in op.inputs]
        if op.has_attribute_holes:
            groups = select_groups(op.definition, inputs, required[op])
            groups = order_groups(groups, op, source, inputs)
            for binding in op.unresolved_axis_bindings:
                P = inputs[binding.operand].axes
                G = groups[binding.group]
                C[binding.attribute] = positions(P, G, M, partitions)
            complete_other_attributes(op, inputs, groups, match, C)

        completed = substitute_attributes(op, C)
        V[op] = lower_and_check(completed, inputs, M, partitions)

    check_equivalence(V[rule.lhs], V[rule.rhs])
    programs = source_only_expressions(C)
    verify_in_fresh_context(rule, programs)
    return programs
\end{lstlisting}
\end{paperalgorithm}

\code{correspond} matches normalized reads, and \code{required_output_axes}
computes $\Keep(t)$ for each replacement intermediate. \code{select_groups}
uses operand membership and reduction roles as in
Section~\ref{sec:keep-axes}. Explicit output-domain groups are separate: they
come from required output axes absent from every operand, and their sizes
determine domain attributes. Ambiguous group assignments remain unresolved.
\code{order_groups} respects fixed bindings and otherwise follows source
and operand orders. \code{positions} abbreviates mapping lookup and position
construction for these ordered groups.

\code{complete_other_attributes} handles supported index coefficients,
offsets, output domains, and output permutations. \code{lower_and_check}
checks each completed operation and publishes its interface and mappings for
consumers. \code{check_equivalence} compares fully lowered values, reduction
domains, output order, shape, and dtype. \code{source_only_expressions}
resolves constructors to source metadata, and the final check reconstructs
the completed rule in a fresh context. The pseudocode omits failure paths:
unsupported or unproved steps return an unresolved result. Successful output
is a program for every missing attribute, ready for match-time evaluation.

\FloatBarrier
\subsection{Correctness and Scope}
\label{sec:correctness}
Synthesis proposes a construction for the replacement; verification must
establish that the construction preserves the original computation. A mapping
query can locate axes without establishing that attributes are evaluable,
operations legal, or computed tensors equal.

\FloatBarrier
\subsubsection{Checking the Completed Rewrite}
After synthesis produces $f$, the verifier substitutes its expressions into
the RHS and checks the completed rule in a fresh context. It reconstructs
tensor meanings and coordinate relationships rather than taking intermediate
synthesis conclusions as established facts. Table~\ref{tab:verification}
lists the checks.
\begin{table}[htbp]
\caption{Checks required for a completed semantic rewrite.}
\label{tab:verification}
\begin{tabularx}{\linewidth}{@{}lY@{}}\toprule
Check & Required fact\\\midrule
Attribute computation & Constructors depend only on $s$, are evaluable whenever the rule applies, and leave no attribute holes.\\
RHS legality & Constructed attributes satisfy operation definitions, including valid position lists, paired sizes, reduction domains, and access bounds.\\
Value equivalence & Normalized source and completed replacement agree in coefficients, indexed reads, and reduction domains.\\
Output agreement & Both sides have the same ordered output axes and shape, with compatible dtype as fixed by the contract.\\\bottomrule
\end{tabularx}
\end{table}
Read matching against a skeleton assists coordinate construction; it does not
prove value equivalence. That check uses the fully lowered RHS with actual
attributes and reduction axes. The verifier can assume source legality and
the declared guard. Any extra condition needed by construction must follow
from those assumptions and cannot be introduced merely to make checking
succeed. Explicit attribute programs undergo the same semantic checking.
The outcome is an accepted completed rule or an explanation of what could
not be established.

\FloatBarrier
\subsubsection{Correctness Guarantee}
\label{sec:guarantee}
The intended guarantee of successful semantic verification is the contract
in Equation~\eqref{eq:contract} for every legal guarded source instance,
rather than only particular shapes or test values. Specifically, for every
legal $s$ satisfying $\operatorname{Guard}(s)$, $f(s)$ should be evaluable,
make the RHS legal, and satisfy
\begin{equation}
 \forall X\text{ compatible with }s.\quad
 \operatorname{LHS}(s,X)=\operatorname{RHS}(s,f(s),X).
\end{equation}
Equality includes ordered output shape and element values, with element types
fixed as in Section~\ref{sec:problem}. It is relative to the selected arithmetic
policy; algebraic reassociation alone does not justify equality under strict
floating-point evaluation.

This is a correctness contract and intended guarantee, rather than a soundness
theorem established in this draft. Its assumptions include correct operation
definitions, sound primitive metadata and index checks, admissible coordinate
frames, and faithful evaluation of partial attribute programs. Checked
arithmetic, coordinate transport, and guard evaluation are implementation
obligations. A fresh context separates checking from construction, but both
paths share semantic infrastructure; this separation alone does not establish
implementation soundness. The current C++ implementation is not claimed to
have been formally proved sound. We present the verification explanation and
its assumptions without introducing a soundness proof.

\FloatBarrier
\subsubsection{Scope and Limitations}
Failure to establish a rule does not establish that the rule is incorrect.
The method is incomplete and accepts semantic rewrites only when supported
reasoning establishes all required conditions. Ambiguous read correspondence,
unsupported coupled or inverse attribute equations, executable index or guard
computations beyond the checker's reasoning fragment, unavailable arithmetic
laws, and bounded normalization or coordinate work can all prevent acceptance.
For example, several source-coordinate assignments for a replacement read
must remain unresolved unless the checker can establish agreement.

Diagnostics should distinguish malformed rules, unsupported or unresolved
reasoning, and exhausted proof-work limits. TEPL does not claim to find every
valid attribute program or the weakest sufficient guard. The intended tensor
equivalence guarantee applies to semantically verified rules. Structural rules
receive structural checks and do not carry that guarantee. Acceptance records
success of the supported checks under stated conditions; rejection as unproved
is not by itself a counterexample.


\FloatBarrier
\FloatBarrier
\section{Implementation}
\label{sec:implementation}
\FloatBarrier
\subsection{High-Level Architecture}
Figure~\ref{fig:architecture} distinguishes TEPL's three inputs: a rule,
a dialect defining operations and their semantics, and a graph to rewrite.
Rules and dialects feed compilation, whereas the graph enters the runtime
engine directly. The implementation separates rule compilation from rewrite
execution. The additional-engine branch in the figure is an extension target,
not an implemented e-graph integration.


\begin{figuretext}{figures/TEPL architecture.pdf}{Rules and dialect semantics compile into reusable programs that a rewrite engine applies to input graphs.}{fig:architecture}{Rule and dialect cards feed the compiler and compiled JSON. A graph card
bypasses compilation and enters a greedy rewrite loop. A dashed adapter
connects compiled rules to a prospective additional engine.}
TEPL is implemented in C++20 and uses ANTLR to parse operation definitions
and rules. Compilation resolves names and checks types, constructs missing
attributes, checks semantic rules, and lowers completed rules into an
executable representation that can be exported as JSON:
\begin{codeblock}[H]
\begin{lstlisting}
Rule compilation
TEPL source
  -> ANTLR parsing
  -> Name resolution and type checking
  -> Attribute synthesis and semantic verification
  -> Executable rewrite program
  -> JSON export
\end{lstlisting}
\end{codeblock}
Type-checked rules may still contain holes; executable rules have completed
attribute programs, guards, and replacement patterns. JSON serializes this
executable representation for loading without repeating compilation or
semantic verification.

The rewrite engine loads and validates the program, searches the input graph
for an LHS match, evaluates its guard, computes replacement attributes, and
constructs and applies the RHS:
\end{figuretext}

\begin{codeblock}[H]
\begin{lstlisting}
Rewrite execution
Compiled JSON + input graph
  -> Load and validate
  -> Find an LHS pattern match
  -> Evaluate the guard
  -> Compute replacement attributes
  -> Construct and apply the RHS
\end{lstlisting}
\end{codeblock}
A failed guard or invalid attribute evaluation skips the match. The engine
rewrites greedily until no rule changes the graph or an execution limit is
reached. As described in Section~\ref{sec:auto}, generated Rust provides
another execution path with the shared runtime; C++ generation is unavailable.

\FloatBarrier
\subsection{Internal Representations}
Symbolic groups and position expressions use immutable typed nodes.
Types distinguish axes, positions, extents, integers, and Boolean conditions.
Gather, concatenation, and inverse permutation remain symbolic until their
inputs are known. Expression nodes can share subexpressions. Structural keys
identify expressions with the same structure, and each proof context interns
these keys to accelerate comparisons. Hashes support lookup, with structural
comparison resolving collisions.

Value nodes represent constants, operand reads, binary addition, binary
multiplication, and reduction. A separate index-expression representation
handles coordinate selection, concatenation, scaling, and offsets for
calculated reads. Normalization produces a vector of terms, each storing
an integer coefficient, a list of reads, and a symbolic reduction sequence,
as in Section~\ref{sec:normalization}.

\FloatBarrier
\subsection{Mapping Storage and Lookup}
Mapping entries store \code{(interface, group, positions)}. Separate indexes
support interface/group lookup and reverse lookup by interface and position
sequence. Partitions refer to entries through stable identifiers and retain
their own position sequences when these differ from the table's preferred
sequence. Overlap records retain the selected subgroup and positions in both
parent groups.

Lookup first checks existing mappings. If no direct mapping is available,
it follows established relationships and composes position sequences.
Concatenated groups combine component mappings. Traversal through registered
origins uses breadth-first search and avoids revisiting an origin. Results
are cached; registering mappings invalidates the mapping cache, and changes
to established facts invalidate the corresponding reasoning caches. Later
queries can therefore exploit newly discovered relationships.

\FloatBarrier
\subsection{Rule Execution}
The engine searches reachable graph nodes for LHS matches, checking operation
identity and operand structure and recording captured nodes, descriptors, and
shape bindings. Repeated tensor captures must refer to the same node.
The evaluator then processes the guard and replacement expressions using
those bindings. Boolean conditions preserve short-circuit evaluation; a
false guard or evaluation failure skips the match.

The builder reuses captured nodes and constructs new RHS operations with
the computed attributes. Construction occurs on a candidate graph, permitting
operation-contract and output-compatibility validation before committing
the replacement. After a successful replacement the engine restarts its
search. It stops when no candidate changes the graph, a configured limit
is reached, or execution is cancelled.

The mechanisms in this section describe the implementation status recorded
in the supplied system outline; they are separate from the intended universal
contract and from unperformed experiments. \todo{Pin the implementation and
runtime revisions and attach build, supported-operation, semantic-checker,
and generated-Rust validation evidence. The paper repository contains no
implementation artifact with which to independently audit these statements.}

\FloatBarrier
\FloatBarrier
\section{Evaluation Plan}
\label{sec:evaluation}
The evaluation will measure the language's authoring benefit, the additional
work removed by inference, and the coverage and cost of its supported
reasoning. This section specifies methodology; no performance, coverage,
reliability, or authoring-effort results have yet been measured for this draft.

\FloatBarrier
\subsection{Research Questions}
\begin{enumerate}
\item \textbf{RQ1---Authoring effort.} How much does TEPL simplify rewrite
authoring, and how much additional attribute logic does inference remove?
What explicit expressions, conditions, and semantic definitions remain?
\item \textbf{RQ2---Coverage.} Which practical rewrites can be executed,
verified, and synthesized, including rewrites requiring newly constructed
rather than copied descriptors?
\item \textbf{RQ3---Scalability.} How do composition depth, partition overlap,
and expression size affect synthesis?
\item \textbf{RQ4---Generality.} Can semantic definitions alone add support
for new operations?
\item \textbf{RQ5---Reliability.} Does checking reject incorrect explicit
constructors, corrupted synthesized constructors, and invalid rewrites?
\end{enumerate}

\FloatBarrier
\subsection{Workloads and Comparisons}
We will select a production-derived corpus with documented inclusion and
exclusion criteria. It will include distribution, transpose/dot transformations,
guarded reassociation, reductions, broadcasts, slices, and supported
convolution rewrites. Each entry will record its production origin and pinned
revision, operation definitions, applicability guard, arithmetic policy,
replacement graph, and whether attributes are reused or reconstructed.
Selection must make unsupported cases visible rather than retaining only
successful examples. \todo{Choose and publish the corpus, revision manifest,
inclusion criteria, exclusions, and rule counts.}

Each rewrite will have three aligned forms: a handwritten host-language
implementation, TEPL with explicit \code{derive} expressions, and TEPL with
inferred attributes. Applicability conditions and numerical policies will
be aligned before comparing code size or coverage. The host-to-explicit-TEPL
comparison isolates the language's contribution, while explicit-to-inferred
TEPL measures the additional work removed by synthesis. Simple descriptor
reuse and complex reconstruction will be reported separately. Production
code will be classified into attribute transformation, profitability, graph
construction, and unrelated infrastructure. The pinned XLA reassociation
example in Section~\ref{sec:introduction} is a starting point, not a
completed size measurement. \todo{Prepare equivalent implementations and
review the source-code classifications for each baseline.}

A feasible bounded-search baseline will assess synthesis cost, with explicit
rank bounds, attribute domains, search budget, and failure classification.
Its finite search space must be stated when comparing it to symbolic
construction. Related verifiers will be compared only on shared expressible
tasks, separating verification coverage from synthesis coverage. Required
user guards and operation-definition effort will accompany rule-length
comparisons. \todo{Select baselines, implement the bounded search, and
identify the shared task subset and numerical assumptions.}

\FloatBarrier
\subsection{Measurements and Ablations}
\paragraph{Authoring effort and semantic-definition cost.}
For each rule, we will measure attribute-code size and count remaining explicit
attribute expressions. We will also record user-supplied guards and the size
of required operation definitions. Initial specification cost and reuse across
rules will both be reported. Corpus-wide authored code is needed alongside
per-rule savings to reveal logic moved into shared builtins or semantic
definitions. Code size is an effort proxy and cannot by itself establish
usability. \todo{Specify counting conventions and collect per-rule and
corpus-wide authored-code measurements, including shared infrastructure.}

\paragraph{Coverage and generality.}
Execution, verification, and synthesis coverage will use separate denominators.
Each attempted task will be classified as accepted, unsupported, invalid,
or budget-limited, with the failed stage and diagnostic recorded. Attribute
reconstruction will be distinguished from copying an existing descriptor.
To address RQ4, we will document the definitions added for selected operations
and any implementation changes required beyond those definitions.
\todo{Populate coverage classifications and the operation-extension study;
report all denominators and failed cases.}

\paragraph{Cost and scalability.}
We will measure synthesis time, verification time, peak memory, emitted
constructor size, and matching-time cost in the greedy engine and generated
Rust. Compilation costs and application costs will be reported separately,
since a reusable constructor can be compiled once and evaluated many times.
Scaling workloads will vary source composition depth, overlap structure,
repeated factors, expression size, and concrete constructor inputs.
\todo{Fix hardware, software versions, resource limits, repetitions, timing
boundaries, and variability reporting; then collect times, memory, constructor
sizes, and runtime costs.}

Ablations will disable reusable interfaces, lazy refinement, and mapping
caches individually to isolate their effects. Each ablation will use the same
rules, budgets, and workload families, recording both successful cost and
changes in completion rate. \todo{Implement the ablations, document their
configuration, and collect results.}

\paragraph{Reliability and boundary cases.}
Independent numerical references will compare source and completed replacement
computations on instantiated metadata. The test matrix will include interleaved
selections, empty groups, multiple contractions, rank-zero tensors, and
supported zero extents. Incorrect explicit constructors and mutations of
synthesized constructors will test whether the checker detects broken position
lists, pairings, accesses, or output layouts. Invalid rewrites and arithmetic
policies that do not permit the needed laws will supply negative cases.
Numerical testing is implementation evidence and supplements universal
reasoning; it does not replace it. \todo{Define independent references,
metadata generation, mutation operators, seeds, expected outcomes, and measured
rejection results.}

The planned result record will associate every rule with authored-code
measurements, guard and semantics costs, the three coverage outcomes,
constructor size, synthesis and verification cost, application cost, and a
failure explanation where applicable. \todo{Populate this record and write
the results and interpretation once measurements exist. All performance
and coverage values are intentionally unfilled.}

\FloatBarrier
\FloatBarrier
\section{Related Work and Novelty Positioning}
\label{sec:related-work}
\paragraph{Arbitrary-rank tensor rewrite verification.}
TensorRight provides automated verification of tensor graph rewrites for
arbitrary rank and size, using aggregated axes to describe unbounded groups
of named axes~\cite{tensorright}. TEPL builds on this grouping idea. Its
design uses ordered group interfaces and position correspondence to construct
replacement-attribute programs for a user-supplied graph. Neither arbitrary-rank
reasoning nor axis grouping alone is a novelty claim. A detailed comparison
must align operator semantics and guards before reporting shared verification
coverage. \todo{Complete the operation-by-operation and task-by-task comparison
with TensorRight.}

\paragraph{Tensor rewrite synthesis and equality saturation.}
TASO generates graph substitutions from operator specifications and verifies
candidate substitutions against supplied operator properties before using
cost-based graph search~\cite{taso}. TENSAT applies equality saturation to
tensor graph superoptimization and uses extraction to select a graph from the
represented alternatives~\cite{tensat}. These systems address generation or
selection of computation graphs. TEPL's synthesis problem, as stated in
Section~\ref{sec:find}, fixes the user-written replacement structure and seeks
reusable functions of source metadata for its attributes. This difference in
problem formulation does not imply that prior systems lack attributes:
TENSAT explicitly represents parameters including strides and axis
permutations~\cite{tensat}. \todo{Audit the parameterization, generated artifacts,
and numerical assumptions of the relevant rules in TASO and TENSAT before
claiming a broader expressiveness distinction.}

\paragraph{Semantics-driven synthesis.}
Souper synthesizes optimizations for LLVM IR using a superoptimization
approach~\cite{souper}. Ruler uses equality saturation to infer rewrite
rules~\cite{ruler}. Their synthesis targets motivate a distinction between
finding a replacement computation or rule set and filling the attributes of
a specified replacement. TEPL targets source-dependent metadata programs and
checks them under operation legality and explicit guards. A systematic
comparison must consider the synthesized language, parameterization,
assumptions, and correctness contract rather than treating all synthesis
systems as interchangeable. \todo{Extend the primary-source review of
superoptimization and rewrite synthesis and substantiate the proposed
rank-polymorphic attribute-program distinction.}

\paragraph{Index notation and affine compiler representations.}
SPORES translates linear algebra into relational algebra and uses a normal
form built from sums of products of indexed atoms~\cite{spores}.
Section~\ref{sec:normalization} uses a related term representation to match
reads and later check completed expressions. The use of this form is not
itself a novelty claim. MLIR's Linalg dialect separates scalar computation
from mappings between iteration and operand spaces; its indexing maps
support structured transformations~\cite{mlir-linalg}. TEPL similarly needs
explicit access correspondence, but uses symbolic arbitrary-length groups
and position expressions to reconstruct attributes of the operations named
in a replacement graph. \todo{Complete comparisons with other index-notation
and affine IR systems, including the limits of calculated-access synthesis.}

\paragraph{Rewrite DSLs and host callbacks.}
egg combines symbolic rewriting with extension mechanisms such as e-class
analyses and custom appliers~\cite{egg,egg-applier}. MLIR's declarative rewrite
rules can capture attributes and invoke \code{NativeCodeCall} for argument
transformation or custom operation building~\cite{mlir-drr}. These interfaces
allow application-specific logic; the attribute-handling burden depends on
what the author supplies. TEPL brings explicit metadata expressions into
the rewrite language and additionally seeks to infer them from operation
semantics. The comparison concerns how transformations are supplied, rather
than whether another framework can represent rich attributes at all.

The candidate distinction to substantiate is compositional synthesis of
reusable, rank-polymorphic attribute programs from user-defined tensor
semantics. The primary sources above support the narrower comparisons made
here, but this draft makes no ``first'' claim. \todo{Complete a systematic
literature review before asserting established novelty.}

\FloatBarrier
\FloatBarrier
\section{Limitations and Future Work}
\label{sec:limitations}
Inference is incomplete. Ambiguous correspondence, unsupported constructions,
and finite proof-work budgets can reject valid rules. Long operation chains
can grow both partition relationships and normalized expressions; the planned
scalability experiments will determine how these growth patterns affect the
implementation.

The semantic scope is restricted for convolutions and broadcasting, and slices
use positive-size, positive-stride semantics. General padding, data-dependent
indexing, coupled attribute equations, and unsupported numerical policies
remain outside the stated scope. A rule may also require metadata computations
or application conditions that the attribute language cannot express.
\todo{Document the exact supported convolution, broadcasting, slice, and
zero-extent fragments against a pinned implementation revision.}

Executable expressions can extend beyond the semantic checker's supported
fragment. Extending runtime evaluation alone therefore does not extend
verification coverage. Future work includes broader expression verification
and a smaller trusted verification kernel. Correct operation definitions,
primitive checks, coordinate transport, checked arithmetic, and guard
execution remain assumptions or obligations at the appropriate layer;
sharing infrastructure between synthesis and verification deserves particular
attention.

Additional execution backends require validation that expression meanings,
evaluation failures, and arithmetic behavior are preserved across paths.
The greedy engine and generated Rust are the execution paths described in
this draft. C++ generation and adapters to additional engines remain future
work; the dashed branch in Figure~\ref{fig:architecture} is an extension point.

\FloatBarrier
\FloatBarrier
\section{Conclusion}
\label{sec:conclusion}
TEPL connects tensor rewrite authoring with attribute-program synthesis.
Users specify rewrite patterns and conditions, then write attribute changes
explicitly or leave them for the compiler to infer. Axes IR describes ordered
coordinate interfaces and values; symbolic mapping composes and refines axis
relationships; normalization supplies read correspondence; and position queries
construct reusable functions of source metadata. The completed replacement
is then checked for attribute evaluability, operation legality, value
equivalence, and output agreement under the declared conditions and arithmetic
policy.

The central result this design seeks to establish is that semantic coordinate
correspondence can generate executable attribute transformations. The draft
presents that mechanism and its intended correctness contract, without claiming
proved implementation soundness, general completeness, or measured practical
benefits. The planned evaluation will test the authoring effort removed by
inference, the semantic-definition work that remains, production-derived
coverage, synthesis cost, and reliability.

\bibliographystyle{ACM-Reference-Format}
\bibliography{references}
\appendix
\FloatBarrier
\FloatBarrier
\section{IR and Artifact Materials}
\label{sec:artifact}
The core value and index-expression forms are listed in
Tables~\ref{tab:value-expressions} and~\ref{tab:index-expressions}. Their
coordinate interpretation uses the mapping invariant in
Equation~\eqref{eq:mapping-invariant}, complete disjoint partitions from
Section~\ref{sec:mapping-algorithm}, and ordered output interfaces. Normalized
terms have the denotation in Equation~\eqref{eq:normalized-term}. Arithmetic
normalization is conditional on the selected laws and on the reduction-domain
and read-independence checks described in Section~\ref{sec:normalization}.
Algorithms~\ref{alg:mapping} and~\ref{alg:synthesis} in the main text give the construction and synthesis
procedures; neither is a claim of completeness.

The accompanying artifact should record each benchmark's definitions, guards,
explicit and synthesized constructors, arithmetic policy, and failure
classification. Reproduction materials should pin the compiler and runtime,
include independent references, and separate compilation, semantic checking,
and application measurements. \todo{Supply the versioned artifact, formal
IR grammar and complete operation denotations, exact algebraic side-condition
catalog, benchmark manifest, build and reproduction commands, environments,
and independent reference tests. These materials are not supplied in this
draft; the correctness contract and assumptions remain as stated in the main
text.}

\end{document}
